\section*{9月25日}
\begin{enumerate}
    \item 在平面直角坐标系 ${xOy}$ 中,已知两点 $S\left( {4,0}\right) ,T\left( {1,0}\right)$ ,动点 $P$ 满足 $\left| {PS}\right|  = 2\left| {PT}\right|$ , 设点 $P$ 的轨迹为 $C$ . 如图,动直线 $l$ 与曲线 $C$ 交于不同的两点 $A,B$ ( $A,B$ 均在 $x$ 轴上方),且 $\angle {ATO} + \angle {BTO} = {180}^{ \circ  }$ .
    \begin{center}
        \includegraphics[width=0.25\linewidth]{bodyxb//src/bx01ch2.2-20.png}
    \end{center}
    
    \begin{enumerate}
        \item 求曲线 $C$ 的方程；
        \item 当 ${A}$ 为曲线 $C$ 与 $y$ 轴正半轴的交点时,求直线 $l$ 的方程;
        \item 是否存在一个定点,使得直线 $l$ 始终经过此定点？若存在,求出定点的坐标；若不存在,请说明理由. 
    \end{enumerate}
    \begin{jieda}
        \begin{enumerate}
            \item 设 ${P}({x}, {y})$ ，由 $|P S|=2|P T|$ 得 $\sqrt{(x-4)^2+y^2}=2 \sqrt{(x-1)^2+y^2}$ ，化简得 $x^2+y^2=4$ ，则曲线 $C$ 的方程为 $x^2+y^2=4$ ；
            \item 由题意知 $A(0,2)$ ，设 ${B}\left({x}_2, {y}_2\right)$ ，
            依题意可知直线 $l$ 的斜率存在，设直线 $l$ 的方程为 $y=k x+2$ ，
            由 $\angle A T O+\angle B T O=180^{\circ}$ ，得 $k_{A T}+k_{B T}=0$ ，
            则 $\left\{\begin{array}{l}-2+\dfrac{y_2}{x_2-1}=0 \\ x_2^2+y_2^2=4\end{array}\right.$ ，所以 $\left\{\begin{array}{l}x_2=0 \\ y_2=-2\end{array}\right.$（舍去）或 $\left\{\begin{array}{l}x_2=\dfrac{8}{5} \\ y_2=\dfrac{6}{5}\end{array}\right.$ ，即 $B\left(\dfrac{8}{5}, \dfrac{6}{5}\right)$ ，
            则 $k=\dfrac{\dfrac{6}{5}-2}{\dfrac{8}{5}-0}=-\dfrac{1}{2}$ ，
            则直线 $l$ 方程为 $y=-\dfrac{1}{2} x+2$ ；
            
            
            \item 设直线 $l$ 方程为 $y=k x+b$ ，设 ${A}\left({x}_1, {y}_1\right), {B}\left({x}_2, {y}_2\right)$ ，
            
            $\because \angle A T O+\angle B T O=180^{\circ}, \therefore k_{A T}+k_{B T}=0$
            
            $\therefore \dfrac{y_1}{x_1-1}+\dfrac{y_2}{x_2-1}=0 \Longrightarrow y_1\left(x_2-1\right)+y_2\left(x_1-1\right)=0$
             
            $\therefore \left(k x_1+b\right)\left(x_2-1\right)+\left(k x_2+b\right)\left(x_1-1\right)=0$
            
            $\therefore 2 k x_1 x_2+(b-k)\left(x_1+x_2\right)-2 b=0$
            
            
            联立方程 $\left\{\begin{array}{l}x^2+y^2=4 \\ y=k x+b\end{array}\right.$ ，得 $\left(k^2+1\right) x^2+2 k b x+b^2-4=0$ ，因此$\left\{\begin{array}{l}
                 x_1+x_2=\dfrac{-2 k b}{k^2+1} \\
                 x_1 x_2=\dfrac{b^2-4}{k^2+1}
            \end{array}\right.$

            $\therefore 2 k \dfrac{b^2-4}{k^2+1}+(b-k)\left(\dfrac{-2 k b}{k^2+1}\right)-2 b=0$，即$2 k\left(b^2-4\right)-2 k b(b-k)-2 b\left(k^2+1\right)=0$

            故$b = -4k$，则直线 $l$ 始终经过此定点 $(4,0)$ ．
            
        \end{enumerate}
    \end{jieda}
    
    \item 已知椭圆 $C : \dfrac{{x}^{2}}{{a}^{2}} + \dfrac{{y}^{2}}{{b}^{2}} = 1\left( {a > b > 0}\right)$ 的上顶点和两焦点构成的三角形为等腰直角三角形,且面积为 2,点 $M$ 为椭圆 $C$ 的右顶点.
    \begin{enumerate}
        \item 求椭圆 $C$ 的方程
        \item 若经过点 $P\left( {t,0}\right)$ 的直线 $l$ 与椭圆 $C$ 交于 $A,B$ 两点,实数 $t$ 取何值时以 ${AB}$ 为直径的圆恒过点 $M$ ?
    \end{enumerate}
    \begin{jieda}
        \begin{enumerate}
            \item 由题意知 $\left\{  {\begin{array}{l} b = c, \\  {bc} = 2, \end{array}\;\text{ 解得 }b = c = \sqrt{2},}\right.$，又 ${a}^{2} - {b}^{2} = {c}^{2}$ ,则 $a = 2$，所以椭圆 $C$ 的方程为 $\dfrac{{x}^{2}}{4} + \dfrac{{y}^{2}}{2} = 1$ .
            \item $M\left( {2,0}\right)$
            \begin{dingautolist}{172}
                \item 若直线 $l$ 的斜率不存在,则直线 $l$ 的方程为 $x = t\left( {-2 < t < 2}\right)$ ,
                
                此时 $A\left( {t,\sqrt{2 - \dfrac{{t}^{2}}{2}}}\right) ,B\left( {t, - \sqrt{2 - \dfrac{{t}^{2}}{2}}}\right)$ ,
                
                由 $\vv{MA} \cdot  \vv{MB} = 0$ 得 $\left( {t - 2,\sqrt{2 - \dfrac{{t}^{2}}{2}}}\right) \cdot \left( {t - 2, - \sqrt{2 - \dfrac{{t}^{2}}{2}}}\right)  = 0$，即$(t-2)^2=2-\dfrac{t^2}{2}$
                
                解得 $t = \dfrac{2}{3}$ 或 $t = 2$ (舍),即 $t = \dfrac{2}{3}$ .
                \item 若直线 $l$ 的斜率存在,不妨设直线 $l : y = k\left( {x - t}\right) ,A\left( {{x}_{1},{y}_{1}}\right) ,B\left( {{x}_{2},{y}_{2}}\right)$ 

                由题意知 $\vv{MA} \cdot  \vv{MB} = 0$ ,即 $\left( {{x}_{1} - 2,{y}_{1}}\right)  \cdot  \left( {{x}_{2} - 2,{y}_{2}}\right)  = 0$，即$\left( {1 + {k}^{2}}\right) {x}_{1}{x}_{2} - \left( {2 + {k}^{2}t}\right) \left( {{x}_{1} + {x}_{2}}\right)  + 4 + {k}^{2}{t}^{2} = 0$
                
                联立 $\left\{  \begin{array}{l} y = k\left( {x - t}\right) , \\  
                \dfrac{x^2}{4}+ \dfrac{{y}^{2}}{2} = 1, \end{array}\right.$，消去$y$得 $\left( {1 + 2{k}^{2}}\right) {x}^{2} - 4{k}^{2}{tx} + 2{k}^{2}{t}^{2} - 4 = 0$ .
                
                根据韦达定理有$\left\{\begin{array}{l}
                     {x}_{1} + {x}_{2} = \dfrac{4{k}^{2}t}{1 + 2{k}^{2}}  \\
                     {x}_{1}{x}_{2} = \dfrac{2{k}^{2}{t}^{2} - 4}{1 + 2{k}^{2}}
                \end{array}\right.$
                
                故有$\left( {1 + {k}^{2}}\right) \left( {2{k}^{2}{t}^{2} - 4}\right)  - \left( {2 + {k}^{2}t}\right)  \cdot  4{k}^{2}t + \left( {4 + {k}^{2}{t}^{2}}\right) \left( {1 + 2{k}^{2}}\right)  = 0$ ,
                
                整理得 ${k}^{2}\left( {3{t}^{2} - {8t} + 4}\right)  = 0$ ,因为 $k$ 不恒为 0,
                
                故解得 $t = \dfrac{2}{3}$ 或 $t = 2$ (舍)
            \end{dingautolist}
            综上，当 $t = \dfrac{2}{3}$ 时，以 ${AB}$ 为直径的圆恒过点 $M$ .
        \end{enumerate}
    \end{jieda}
    
    \item 椭圆 $\dfrac{x^2}{4} + \dfrac{{y}^{2}}{3} = 1$ 的上有两点 $A\text{ 、 }B$ ,点 $B$ 在直线 $x = 4$ 上的射影是点 ${B}^{\prime }$ ,若直线 ${AB}$ 过右焦点,则直线 $A{B}^{\prime }$ 必定经过的定点坐标为\blkda{$\left(\dfrac{5}{2}, 0\right)$}.
    \begin{jieda}
        设$A(x_1, y_1), B(x_2, y_2)， B'\left(4, y_2\right)$，右焦点$F(1, 0)$

        设$AB: y = k(x-1), k = \dfrac{y_2 - y_1}{x_2 - x_1}$，设$AB': y - y_1 = k'(x - x_1), k' = \dfrac{y_2-y_1}{4 - x_1} = \dfrac{x_2-x_1}{4 - x}k$，故$AB'$可以表示为$y - k(x_1-1) = \dfrac{x_2-x_1}{4 - x_1}k(x-x_1)$

        易知$x_1 \neq 4$，故欲证$AB'$过定点，可以先去分母，即证方程$\left(4 - x_1\right)y + k(x_1-1)\left(x_1-4\right) = (x_2 - x_1)kx -x_1x_2k + kx_1^2$存在一组对任意的参数恒成立的解$(x_0, y_0)$.

        易知该方程对于$x_1, x_2$应是轮换对称的，故可以将其整理成$x_1, x_2$为主元的式子，即$\left(-y -5k+kx\right)x_1  - (kx)x_2 + \left(4y+4k+x_1x_2k\right) = 0$. 最后一项对于$x_1, x_2$是轮换对称的，因此必有前两项的系数相等，即$y + 5k - kx = kx$

        若$AB'$过定点$(x_0, y_0)$，则对任意的$x_1, x_2$，等式代入$(x_0, y_0)$均成立，则应有$y_0 + 5k - kx_0 = kx_0$恒成立，而$k$并非定值，因此$x_0 = \dfrac{5}{2}, y_0 = 0$可以使得等式恒成立，故若存在满足题干要求的顶点，则必为$\left(\dfrac{5}{2}, 0\right)$

        下证$A{B}^{\prime }$ 必定经过的定点$\left(\dfrac{5}{2}, 0\right)$：

        将$\left(\dfrac{5}{2}, 0\right)$代入$\left(-y -5k+kx\right)x_1  - (kx)x_2 + \left(4y+4k+x_1x_2k\right) = 0$，可得$LHS = -\dfrac{5}{2}k(x_1+x_2)+(4+x_1x_2)k = k\left(4+x_1x_2-\dfrac{5}{2}(x_1+x_2)\right)$
        
        联立$\left\{\begin{array}{l}
             y = k(x-1) \\
             \dfrac{x^2}{4} + \dfrac{{y}^{2}}{3} = 1
        \end{array}\right.$，消去$y$可得$3x^2+4k^2(x-1)^2 = 12$，整理得$(3+4k^2)x^2-8k^2x+(4k^2-12) = 0$，根据韦达定理有$\left\{\begin{array}{l}
             x_1+x_2 = \dfrac{8k^2}{4k^2+3} \\
             x_1x_2 = \dfrac{4k^2-12}{4k^2+3}
        \end{array}\right.$，代入可得$LHS = k\left(4+x_1x_2-\dfrac{5}{2}(x_1+x_2)\right) = 0 = RHS$
    \end{jieda}
\end{enumerate}
\section*{10月9日}
\begin{enumerate}
    \item 椭圆 $C : \dfrac{{x}^{2}}{4} + {y}^{2} = 1$ 的左、右焦点分别为 ${F}_{1},{F}_{2},O$ 为坐标原点,以下四个命题中正确的是\blkda{\ding{172}\ding{173}\ding{175}}
        \begin{dingautolist}{172}
            \item 若过点 ${F}_{2}$ 的直线与椭圆 $C$ 交于 $A,B$ 两点,则 $\bigtriangleup  {AB}{F}_{1}$ 的周长为 8
            \item 椭圆 $C$ 上存在点 $P$ ,使得 $\vv{P{F}_{1}} \cdot  \vv{P{F}_{2}} = 0$
            \item 椭圆 $C$ 的离心率为 $\dfrac{1}{2}$
            \item 若 $P$ 为椭圆 $\dfrac{{x}^{2}}{4} + {y}^{2} = 1$ 上一点, $Q$ 为圆 ${x}^{2} + {y}^{2} = 1$ 上一点,则点 $P,Q$ 的最大距离为 3
        \end{dingautolist}
        \begin{jieda}
            \begin{dingautolist}{172}
                \item 根据定义可知该三角形周长为$4a = 8$
                \item 根据焦点三角形结论可知$\angle F_1PF_2$的取值范围是$\left(0, \dfrac{2\pi}{3}\right]$
                \item $a = 2, b = 1, c = \sqrt{3}, e = \dfrac{\sqrt{3}}{2}$
                \item 圆和椭圆内切，结合图像可以判定命题正确.
            \end{dingautolist}
        \end{jieda}
    \item 已知椭圆 $C : \dfrac{{x}^{2}}{4} + \dfrac{{y}^{2}}{3} = 1$ 的左、右焦点分别是 ${F}_{1},{F}_{2},M\left( {\dfrac{4}{3},{y}_{0}}\right)$ 为椭圆 $C$ 上一点,则下列结论正确的是\blkda{\ding{172}\ding{174}}
        \begin{dingautolist}{172}
            \item $\bigtriangleup M{F}_{1}{F}_{2}$ 的周长为 6
            \item $\bigtriangleup M{F}_{1}{F}_{2}$ 的面积为 $\dfrac{\sqrt{15}}{2}$
            \item $\bigtriangleup M{F}_{1}{F}_{2}$ 的内切圆的半径为 $\dfrac{\sqrt{15}}{9}$
            \item $\bigtriangleup M{F}_{1}{F}_{2}$ 的外接圆的直径为 $\dfrac{32}{11}$
        \end{dingautolist}
        \begin{jieda}
            不妨设$M$在第一象限，则$M\left(\dfrac{4}{3}, \dfrac{\sqrt{15}}{3}\right)$
            \begin{dingautolist}{172}
               \item 根据定义，三角形周长为$2a+2c = 4+2 = 6$
               \item $S = \dfrac{1}{2} \times 2 \times \dfrac{\sqrt{15}}{3} = \dfrac{\sqrt{15}}{3}$
               \item $S = \dfrac{1}{2} \times 6 \times r = \dfrac{\sqrt{15}}{3}$，因此$r = \dfrac{\sqrt{15}}{9}$
               \item $\cos \angle {F}_{1}M{F}_{2} = \dfrac{\dfrac{64}{9} + \dfrac{16}{9} - 4}{2 \times  \dfrac{8}{3} \times  \dfrac{4}{3}} = \dfrac{11}{16}$,则 $\sin \angle {F}_{1}M{F}_{2} = \sqrt{1 - {\cos }^{2}\angle {F}_{1}M{F}_{2}} = \dfrac{3\sqrt{15}}{16}$\\
               根据正弦定理 $\bigtriangleup M{F}_{1}{F}_{2}$ 的外接圆的直径为 $\dfrac{\left| {F}_{1}{F}_{2}\right| }{\sin \angle {F}_{1}M{F}_{2}} = \dfrac{2}{\dfrac{3\sqrt{15}}{16}} = \dfrac{{32}\sqrt{15}}{45}$.
               
               
               \ps 根据焦点三角形面积公式计算$tan\left(\dfrac{1}{2}\angle {F}_{1}M{F}_{2}\right) = \dfrac{\sqrt{15}}{9}$，然后利用万能公式$sin\angle {F}_{1}M{F}_{2} = \dfrac{2tan\left(\dfrac{1}{2}\angle {F}_{1}M{F}_{2}\right)}{1+tan^2\left(\dfrac{1}{2}\angle {F}_{1}M{F}_{2}\right)} = \dfrac{3\sqrt{15}}{16}$
            \end{dingautolist}
        \end{jieda}
    \item 已知 $O$ 为坐标原点,椭圆 $C : \dfrac{{x}^{2}}{{a}^{2}} + \dfrac{{y}^{2}}{{b}^{2}} = 1\left( {a > b > 0}\right)$ 的离心率为 $\dfrac{\sqrt{6}}{3}$ ,且经过点 $P\left( {\sqrt{6},1}\right)$ .
    \begin{enumerate}
        \item 求椭圆 $C$ 的方程
        \item 直线 $l$ 与椭圆 $C$ 交于 $A,B$ 两点,直线 ${OA}$ 的斜率为 ${k}_{1}$ ,直线 ${OB}$ 的斜率为 ${k}_{2}$ ,且 ${k}_{1}{k}_{2} =  - \dfrac{1}{3}$ ,求 $\vv{OA} \cdot  \vv{OB}$ 的取值范围.
    \end{enumerate}
    \begin{jieda}
        \begin{enumerate}
            \item 由题意可得 $\left\{  \begin{array}{l} \dfrac{c}{a} = \dfrac{\sqrt{6}}{3}, \\  \dfrac{6}{{a}^{2}} + \dfrac{1}{{b}^{2}} = 1, \end{array}\right.$，因为${a}^{2} = {b}^{2} + {c}^{2}$，解得 $a = 3,b = \sqrt{3}$ .
            
            所以椭圆 $C$ 的方程为 $\dfrac{{x}^{2}}{9} + \dfrac{{y}^{2}}{3} = 1$ .
            \item 
            \begin{enumerate}
                \item 方法1

                联立$OA$和椭圆方程$\left\{\begin{array}{l}
                     y = k_1x  \\
                     \dfrac{{x}^{2}}{9} + \dfrac{{y}^{2}}{3} = 1 
                \end{array}\right.$，消去$y$可得$x_A^2 = \dfrac{9}{1+3k_1^2}$，考虑$x_A = \dfrac{3}{\sqrt{1+3k_1^2}}$的情况，另一种情况数量积为相反数.
                
                若$k_1>0$，可得$x_B = \dfrac{3}{\sqrt{1+3k_2^2}} = \dfrac{3\sqrt{3}k_1}{\sqrt{1+3k_1^2}}$

                故$A\left(\dfrac{3}{\sqrt{1+3k_1^2}}, \dfrac{3k_1}{\sqrt{1+3k_1^2}}\right)$，$B\left(\dfrac{3\sqrt{3}k_1}{\sqrt{1+3k_1^2}}, \dfrac{-\sqrt{3}}{\sqrt{1+3k_1^2}}\right)$

                故$\vv{OA} \cdot \vv{OB} = \dfrac{6\sqrt{3}k_1}{1+3k_1^2} = \dfrac{6\sqrt{3}}{3k_1+\dfrac{1}{k_1}} \in (0, 3]$

                根据对称性可知$k_1 < 0$时，范围相同

                故根据对称性可知$x_A < 0$时$\vv{OA} \cdot \vv{OB} \in \left\lbrack -3,0\right)$，故$\vv{OA} \cdot \vv{OB} \in \left\lbrack  {-3,0)\cup (0,3}\right\rbrack$

                \item 方法2

                设 $A\left( {{x}_{1},{y}_{1}}\right) ,B\left( {{x}_{2},{y}_{2}}\right)$
                \begin{dingautolist}{172}
                    \item 当直线 $l$ 的斜率不存在时, ${x}_{2} = {x}_{1},{y}_{2} =  - {y}_{1}$ ,则 ${k}_{1}{k}_{2} =  - \dfrac{{y}_{1}^{2}}{{x}_{1}^{2}} =  - \dfrac{1}{3}$ ,又 $\dfrac{{x}_{1}^{2}}{9} + \dfrac{{y}_{1}^{2}}{3} = 1$ ,解得 ${x}_{1}^{2} = \dfrac{9}{2}$ ,则 $\vv{OA} \cdot  \vv{OB} = {x}_{1}^{2} - {y}_{1}^{2} = \dfrac{2}{3}{x}_{1}^{2} = 3$ .
                    \item 当直线 $l$ 的斜率存在时,设 $l : y = {kx} + t$，联立直线和椭圆方程有 $\left\{  \begin{array}{l} y = {kx} + t, \\  \dfrac{x^2}{9} + \dfrac{{y}^{2}}{3} = 1, \end{array}\right.$
                    
                    消去 $y$ 得 $\left( {1 + 3{k}^{2}}\right) {x}^{2} + {6ktx} + 3{t}^{2} - 9 = 0$，根据题意可得$\Delta  = {12}\left( {3 + 9{k}^{2} - {t}^{2}}\right)  > 0$ ,
                    
                    根据韦达定理有$\left\{  {\begin{array}{l} {x}_{1} + {x}_{2} = \dfrac{-{6kt}}{1 + 3{k}^{2}}, \\  {x}_{1}{x}_{2} = \dfrac{3{t}^{2} - 9}{1 + 3{k}^{2}}, \end{array}}\right.$
                    
                    ${k}_{1}{k}_{2} = \dfrac{{y}_{1}{y}_{2}}{{x}_{1}{x}_{2}} =  - \dfrac{1}{3}$，代入直线方程可得$\dfrac{{y}_{1}{y}_{2}}{{x}_{1}{x}_{2}} = \dfrac{\left( {k{x}_{1} + t}\right) \left( {k{x}_{2} + t}\right) }{{x}_{1}{x}_{2}} = {k}^{2} + \dfrac{{kt}\left( {{x}_{1} + {x}_{2}}\right)  + {t}^{2}}{{x}_{1}{x}_{2}} = {k}^{2} + \dfrac{\dfrac{-6{k}^{2}{t}^{2}}{1 + 3{k}^{2}} + {t}^{2}}{\dfrac{3{t}^{2} - 9}{1 + 3{k}^{2}}} = \dfrac{{t}^{2} - 9{k}^{2}}{3{t}^{2} - 9} =  - \dfrac{1}{3},$
                    
                    整理得 $2{t}^{2} = 9{k}^{2} + 3$且$t^2 \neq 3$，故${t}^{2} \geq \dfrac{3}{2}$且$t^2 \neq 3$，经检验$\Delta  > 0$ 恒成立.
                    
                    $\vv{OA} \cdot  \vv{OB} = {x}_{1}{x}_{2} + {y}_{1}{y}_{2} = {x}_{1}{x}_{2} - \dfrac{1}{3}{x}_{1}{x}_{2} = \dfrac{2}{3}{x}_{1}{x}_{2} = \dfrac{2}{3} \cdot  \dfrac{3{t}^{2} - 9}{1 + 3{k}^{2}} = 3 \cdot  \dfrac{{t}^{2} - 3}{{t}^{2}} = 3\left( {1 - \dfrac{3}{{t}^{2}}}\right) \in  \lbrack  - 3,0) \cup  \left( {0,3}\right)$ .
                \end{dingautolist}
                综上, $\vv{OA} \cdot  \vv{OB}$ 的取值范围为 $\left\lbrack  {-3,0)\cup (0,3}\right\rbrack$ .
            \end{enumerate}
        \end{enumerate}
    \end{jieda}
    \item 已知 ${F}_{1}$ 、 ${F}_{2}$ 分别为椭圆 $W : \dfrac{{x}^{2}}{4} + {y}^{2} = 1$ 的左、右焦点, $M$ 为椭圆 $W$ 上的一点.
    \begin{enumerate}
        \item 若点 $M$ 的坐标为 $\left( {1,m}\right) \left( {m > 0}\right)$ ,求 $\bigtriangleup {F}_{1}M{F}_{2}$ 的面积；
        \item 若点 $M$ 的坐标为 $\left( {{x}_{0},{y}_{0}}\right)$ ,且 $\angle {F}_{1}M{F}_{2}$ 是钝角,求横坐标 ${x}_{0}$ 的范围；
        \item 若点 $M$ 的坐标为$(0,1)$,且直线 $y = {kx} - \dfrac{3}{5}\left( {k \in  \mathbf{R}}\right)$ 与椭圆 $W$ 交于不同两点 $A\text{ 、 }B$ , 求证: $\vv{MA} \cdot  \vv{MB}$ 为定值,并求出该定值.
    \end{enumerate}
    \begin{jieda}
        \begin{enumerate}
            \item 将$M$坐标带入椭圆方程，解得$m = \dfrac{\sqrt{3}}{2}$，因此$S_{\bigtriangleup {F}_{1}M{F}_{2}} = \dfrac{1}{2} \times (2c) \times (m) = \dfrac{1}{2} \times (2\sqrt{3}) \times \dfrac{\sqrt{3}}{2} = \dfrac{3}{2}$
            \item 设$\theta = \angle{F_1MF_2}$根据焦点三角形的面积结论可知$\dfrac{1}{2} \times (2\sqrt{3}) \times y_0 = S = tan\dfrac{\theta}{2}$，因为$\theta$是钝角，所以$tan\dfrac{\theta}{2} > 1$，因此$S>1$，当$M$为上顶点时$S$取最大值$\sqrt{3}$，因此$S \in (1, \sqrt{3})$，即$y_0 \in \left(\dfrac{\sqrt{3}}{3}, 1\right)$，根据椭圆方程可知$x_0 \in \left(-\dfrac{2\sqrt{6}}{3}, \dfrac{2\sqrt{6}}{3}\right)$
            \item 易知直线过定点$N\left(0, -\dfrac{3}{5}\right)$在椭圆内部，因此必有两个交点. 设 $A\left( {{x}_{1},{y}_{1}}\right) ,B\left( {{x}_{2},{y}_{2}}\right)$ 则$\vv{MA} = (x_1, y_1-1), \vv{MB} = (x_2, y_2-1)$
            \begin{enumerate}
                \item 方法1
                
                $\vv{MA} \cdot \vv{MB} = x_1x_2+(y_1-1)(y_2-1) = x_1x_2+\left(kx_1-\dfrac{8}{5}\right)\left(kx_2-\dfrac{8}{5}\right) = (1+k^2)x_1x_2-\dfrac{8k}{5}(x_1+x_2)+\dfrac{64}{25}$

                联立方程组 $\left\{  \begin{array}{l} \dfrac{{x}^{2}}{4} + {y}^{2} = 1, \\  y = {kx} - \dfrac{3}{5}, \end{array}\right.$ 消去$y$有$x^2+4\left(kx-\dfrac{3}{5}\right) = 4$，整理得$(4k^2+1)x^2-\dfrac{24}{5}kx-\dfrac{64}{25} = 0$

                因此根据韦达定理有$\left\{\begin{array}{l}
                     x_1+x_2 = \dfrac{24k}{5(4k^2+1)} \\
                     x_1x_2 = -\dfrac{64}{25(4k^2+1)}
                \end{array}\right.$，代入可得$\vv{MA} \cdot \vv{MB} =  (1+k^2)x_1x_2-\dfrac{8k}{5}(x_1+x_2)+\dfrac{64}{25} = (1+k^2)\left(-\dfrac{64}{25(4k^2+1)}\right) - \dfrac{8k}{5}\left(\dfrac{24k}{5(4k^2+1)}\right) + \dfrac{64}{25} = 0$
                
            
                \item 方法2
                
                $\vv{MA} \cdot \vv{MB} = x_1x_2 + y_1y_2-(y_1+y_2)+1 = \dfrac{1}{k^2}\left(y_1+\dfrac{3}{5}\right)\left(y_2+\dfrac{3}{5}\right) + y_1y_2-(y_1+y_2)+1 = \dfrac{1+k^2}{k^2}y_1y_2 + \left(\dfrac{3}{5k^2}-1\right)(y_1+y_2)+\dfrac{9+25k^2}{25k^2}$\\
                联立方程组 $\left\{  \begin{array}{l} \dfrac{{x}^{2}}{4} + {y}^{2} = 1, \\  y = {kx} - \dfrac{3}{5}, \end{array}\right.$ 消去$x$有$\dfrac{1}{4k^2}\left(y+\dfrac{3}{5}\right)^2+y^2 = 1$，整理得$\dfrac{1+4k^2}{4k^2}y^2+\dfrac{3}{10k^2}y+\dfrac{9-100k^2}{100k^2} = 0$，即$(25+100k^2)y^2+30y+(9-100k^2) = 0$\\
                因此根据韦达定理有$\left\{\begin{array}{l}
                     y_1+y_2 = \dfrac{-30}{25+100k^2} \\
                     y_1y_2 = \dfrac{9-100k^2}{25+100k^2}
                \end{array}\right.$，代入可得$\vv{MA} \cdot \vv{MB} = \dfrac{1+k^2}{k^2} \cdot \dfrac{9-100k^2}{25+100k^2} + \left(\dfrac{3}{5k^2}-1\right)\cdot \dfrac{-30}{25+100k^2} +\dfrac{9+25k^2}{25k^2} = \dfrac{(1+k^2)(9-100k^2) -6(3-5k^2)+(1+4k^2)(9+25k^2)}{k^2(25+100k^2)} = 0$
            \end{enumerate}
            
        \end{enumerate}
    \end{jieda}

    \item 设 ${F}_{1}\text{ 、 }{F}_{2}$ 分别是椭圆 $\dfrac{{x}^{2}}{4} + {y}^{2} = 1$ 的左、右焦点.
    \begin{enumerate}
        \item 若 $P$ 是该椭圆上的一个动点,求 $\vv{P{F}_{1}} \cdot  \vv{P{F}_{2}}$ 的最大值和最小值;
        \item 设过定点 $M\left( {0,2}\right)$ 的直线 $l$ 与椭圆交于不同的两点 $A$ 、 $B$ ,且 $\angle {AOB}$ 为锐角(其中 $O$ 为坐标原点),求直线 $l$ 的斜率 $k$ 的取值范围.
    \end{enumerate}
    \begin{jieda}
        \begin{enumerate}
            \item 根据题意可知$a = 2, b = 1, c = \sqrt{3}$，因此$F_1(-\sqrt{3}, 0), F_2(\sqrt{3}, 0)$，设$P(x_0, y_0)$，则$\vv{P{F}_{1}} = (-\sqrt{3}-x_0, -y_0), \vv{PF_2} = (\sqrt{3}-x_0, -y_0)$，因此$\vv{P{F}_{1}} \cdot \vv{P{F}_{2}} = (x_0^2-3)+y_0^2$，代入椭圆方程有$\vv{P{F}_{1}} \cdot \vv{P{F}_{2}} = (x_0^2-3)+1-\dfrac{x_0^2}{4} = \dfrac{3x_0^2}{4}-2 \in [-2, 1]$，因此最小值为$-2$，最大值为$1$
            \item 设$y = kx+2$，$A(x_1, y_2), B(x_2, y_2)$ \\
            联立直线和椭圆方程$\left\{\begin{array}{l}
                 y = kx+2 \\
                 \dfrac{{x}^{2}}{4} + {y}^{2} = 1
            \end{array}\right.$，消去$y$，得到$\dfrac{x^2}{4} + (kx+2)^2 = 1$，整理得$\left(k^2+\dfrac{1}{4}\right)x^2+4kx+3=0$，$\Delta = 16k^2-12\left(k^2+\dfrac{1}{4}\right) = 4k^2-3$，因此当$k \in \left(-\infty, -\dfrac{\sqrt{3}}{2}\right) \cup \left(\dfrac{\sqrt{3}}{2}, +\infty\right)$时存在两个交点. \\
            在此基础上，$\angle AOB$为锐角即$\vv{OA} \cdot \vv{OB} > 0$，因此有$x_1x_2+y_1y_2 > 0$，代入直线方程有$x_1x_2+(kx_1+2)(kx_2+2) > 0$，化简得$(k^2+1)x_1x_2+2k(x_1+x_2)+4>0$，根据韦达定理有$\left\{\begin{array}{l}
                 x_1+x_2 = \dfrac{-4k}{k^2+\dfrac{1}{4}} = \dfrac{-16k}{1+4k^2}  \\
                 x_1x_2 = \dfrac{3}{k^2+\dfrac{1}{4}} = \dfrac{12}{1+4k^2}
            \end{array}\right.$，代入不等关系可得$(k^2+1)\cdot \dfrac{12}{1+4k^2} + 2k \cdot \dfrac{-16k}{1+4k^2} + 4>0$，即$\dfrac{-4k^2+16}{1+4k^2} > 0$，即$k^2 < 4$，因此有$k \in \left(-2, -\dfrac{\sqrt{3}}{2}\right) \cup \left(\dfrac{\sqrt{3}}{2}, 2\right)$
        \end{enumerate}
    \end{jieda}
\end{enumerate}
\section*{10月23日}
\begin{enumerate}
    \item 过椭圆 $\dfrac{{x}^{2}}{2} + {y}^{2} = 1$ 的左焦点作倾斜角为 ${60}^{ \circ  }$ 的直线,直线与椭圆交于 $A,B$ 两点,则 $\left| {AB}\right|  =$ \blkda{$\dfrac{8\sqrt{2}}{7}$}.
    \begin{jieda}
        $e = \dfrac{\sqrt{2}}{2}$，$p = \dfrac{b^2}{c} = 1$，$|AB| = \dfrac{2ep}{1-e^2cos^2\ang{60}} = \dfrac{\sqrt{2}}{1 - \dfrac{1}{2} \times \dfrac{1}{4}} = \dfrac{8\sqrt{2}}{7}$
    \end{jieda}
    
    \item 过双曲线 ${x}^{2} - {y}^{2} = 4$ 的右焦点 $F$ 作倾斜角为 ${30}^{ \circ  }$ 的直线,交双曲线于 $A,B$ 两点,则弦长 $\left| {AB}\right|  =$ \blkda{8}.
    \begin{jieda}
        $e = \sqrt{2}$，$p = \dfrac{b^2}{c} = \dfrac{4}{2\sqrt{2}} = \sqrt{2}$，$|AB| = \left|\dfrac{2ep}{1-e^2cos^2\ang{30}}\right| = \left|\dfrac{4}{1-2\times \dfrac{3}{4}}\right| = 8$
    \end{jieda}
    
    \item 过椭圆 $3{x}^{2} + 4{y}^{2} = {48}$ 的左焦点 $F$ 引直线 $l$ 交椭圆于 $A$ 、 $B$ 两点,若弦 ${AB}$ 的长为 $\dfrac{48}{7}$ ,则直线 $l$ 的斜率为\blkda{$\pm 1$}.
    \begin{jieda}
        $3{x}^{2} + 4{y}^{2} = {48} \Longleftrightarrow \dfrac{x^2}{16} + \dfrac{y^2}{12} = 1$，$e = \dfrac{1}{2}$，$p = \dfrac{b^2}{c} = 6$，$\dfrac{48}{7} = \dfrac{2ep}{1-e^2cos^2\alpha} = \dfrac{6}{1-\dfrac{cos^2\alpha}{4}}$，故$cos^2\alpha = \dfrac{1}{2}$，故$k = \pm 1$
    \end{jieda}
    
    \item 已知椭圆两顶点 $A\left( {-1,0}\right)$ , $B\left( {1,0}\right)$ ,过焦点 $F\left( {0,1}\right)$ 的直线 $l$ 与椭圆交于 $C,D$ 两点,当 $\left| {CD}\right|  = \dfrac{3\sqrt{2}}{2}$ 时,直线 $l$ 的方程为\blkda[6]{$\sqrt{2}x - y + 1 = 0$ 或 $\sqrt{2}x + y - 1 = 0$}.
    \begin{jieda}
        根据题意可知椭圆的焦点在$y$轴，$e = \dfrac{\sqrt{2}}{2}, p = \dfrac{b^2}{c} = 1$，则$\dfrac{3\sqrt{2}}{2} = \dfrac{2ep}{1-e^2sin^2\alpha} = \dfrac{\sqrt{2}}{1-\dfrac{sin^2\alpha}{2}}$，故$sin^2\alpha = \dfrac{2}{3}$，故$k = \pm \sqrt{2}$，故直线方程为$y = \pm \sqrt{2}x + 1$，即$\sqrt{2}x - y + 1 = 0$ 或 $\sqrt{2}x + y - 1 = 0$.
    \end{jieda}
    
    \item 过双曲线 $C : \dfrac{{x}^{2}}{4} - \dfrac{{y}^{2}}{2} = 1$ 的右焦点 $F$ 的直线 $l$ 与双曲线 $C$ 交于 $A,B$ 两点,若 $\left| {AF}\right|  = \dfrac{2}{3}$ ,则 $\left| {BF}\right|  =$ \blkda{2}.
    \begin{jieda}
            
        \begin{center}
        \includegraphics[width=0.2\textwidth]{bodyxb/src/bo_d3rh8t3ef24c73d0uivg_4_239_1616_311_253_0.jpg}
        \end{center}
        $e = \dfrac{\sqrt{6}}{2}$，$p = \dfrac{b^2}{c} = \dfrac{\sqrt{6}}{3}$

        \begin{enumerate}
            \item 方法1
            $\left|\dfrac{ep}{1-ecos\alpha}\right| = \dfrac{2}{3}$，代入可得$\left|1-\dfrac{\sqrt{6}}{2}cos\alpha\right| = \dfrac{3}{2}$，故$cos\alpha = -\dfrac{\sqrt{6}}{6}$或$\dfrac{5}{6}\sqrt{6}$（舍），故$|BF| = \left|\dfrac{ep}{1+ecos\alpha}\right| = 2$
            
            \item 方法2
            
            因为$\left| {AF}\right|  = \dfrac{2}{3}$，所以$A$在右一支上，故$|AF| = ex_A-a = \dfrac{\sqrt{6}}{2} x_A - 2 = \dfrac{2}{3}$，故$x_A = \dfrac{8\sqrt{6}}{9}$，故$cos\alpha = \dfrac{\sqrt{6} - \dfrac{8\sqrt{6}}{9}}{\dfrac{2}{3}} = \dfrac{\sqrt{6}}{6}$，故$|AB| = \left|\dfrac{2ep}{1-e^2cos^2\alpha}\right| = \dfrac{2}{1-\dfrac{1}{4}} = \dfrac{8}{3}$，故$|BF| = \dfrac{8}{3} - \dfrac{2}{3} = 2.$
        \end{enumerate}
    \end{jieda}
    
    \item 已知椭圆 $C : \dfrac{{x}^{2}}{4} + \dfrac{{y}^{2}}{3} = 1$ ,过右焦点的直线交 $C$ 于 $A,B$ 两点,则 $\left| {AB}\right|$ 的最小值为\blkda{3}.
    \begin{jieda}
        $\left| {AB}\right| = \dfrac{2ep}{1-e^2cos^2\alpha} \leq 2ep = \dfrac{2b^2}{a}$(通径长度为$\dfrac{2b^2}{a}$)
    \end{jieda}
    
    \item 已知椭圆 $C : \dfrac{{x}^{2}}{{a}^{2}} + \dfrac{{y}^{2}}{{b}^{2}} = 1\left( {a > b > 0}\right)$ ,过椭圆 $C$ 的焦点 $F\left( {-\sqrt{2},0}\right)$ 的直线与椭圆相交于 $A$ 、 $B$ 两点. 当线段 ${AB}$ 取得最小值为 2 时,这时椭圆的离心率是\blkda{$\dfrac{\sqrt{2}}{2}$}.
    \begin{jieda}
        $\dfrac{2b^2}{a} = 2$，故$a = b^2$，又因为$c = \sqrt{2}$

        ${a}^{2} - {c}^{2} = a \Longrightarrow  {a}^{2} - 2 = a \Longrightarrow  \left( {a - 2}\right) \left( {a + 1}\right)  = 0 \Longrightarrow  a = 2$ 或 $a =  - 1$ (舍)

        故$e = \dfrac{\sqrt{2}}{2}$
    \end{jieda}
    
    \item 已知椭圆 $C : \dfrac{{x}^{2}}{{a}^{2}} + \dfrac{{y}^{2}}{{b}^{2}} = 1\left( {a > b > 0}\right)$ 的右焦点为 $F\left( {1,0}\right)$ ,过点 $F$ 且垂直于 $x$ 轴的直线与 $C$ 交于 $D,E$ 两点, $O$ 为坐标原点,若 ${OD}\bot {OE}$ ,则 $a =$ \blkda{$\dfrac{1+\sqrt{5}}{2}$}.
    \begin{jieda}
        根据题意可知$\dfrac{b^2}{a} = c = 1$

        
        ${a}^{2} - {c}^{2} = a \Longrightarrow  {a}^{2} - 1 = a \Longrightarrow  a = \dfrac{1+\sqrt{5}}{2}$ 或 $a =  \dfrac{1-\sqrt{5}}{2}$ (舍)
    \end{jieda}
\end{enumerate}
\section*{10月30日}
\begin{enumerate}
    \item 直线 ${l}_{1} : y = x + a$ 和 ${l}_{2} : y = x + b$ 将单位圆 $C : {x}^{2} + {y}^{2} = 1$ 分成长度相等的四段弧,则 ${a}^{2} +$  ${b}^{2} =$ \blkda{2}.
    \begin{jieda}
        数形结合可知$a = 1$，$b = -1$或对称情况.
    \end{jieda}
    \item 若 $P$ 是双曲线 $C : \dfrac{{x}^{2}}{{a}^{2}} - \dfrac{{y}^{2}}{{b}^{2}} = 1\left( {a > 0,b > 0}\right)$ 的渐近线上任意一点,下列正确的是\dotfill{(\kh{C})}
    \xx{存在过点 $P$ 的直线与该双曲线相切}{不存在过点 $P$ 的直线与该双曲线相切}{至多存在一条过点 $P$ 的直线与该双曲线相切}{至多存在一条过点 $P$ 的直线与该双曲线只有一个交点}
    \begin{jieda}
        当$P$为原点的时候，不存在切线，否则至多存在一条切线
    \end{jieda}
    \item 动圆 $C$ 经过点 $F\left( {1,0}\right)$ ,并且与直线 $x =  - 1$ 相切,若动圆 $C$ 与直线 $y = x + 2\sqrt{2} + 1$ 总有公共点,则圆 $C$ 的面积\dotfill{(\kh{D})}
    \xx{有最大值 ${8\pi }$}{有最小值 ${2\pi }$}{有最小值 ${3\pi }$}{有最小值 ${4\pi }$}
    \begin{jieda}
        根据题意易知动圆 $C$的圆心位于抛物线$y^2 = 4x$上，设圆心为$\left(\dfrac{b^2}{4}, b\right)$
        
        若动圆 $C$ 与直线 $y = x + 2\sqrt{2} + 1$ 总有公共点，则圆心到直线的距离应大于等于动圆 $C$ 的半径，即$\dfrac{\left|\dfrac{b^2}{4} - b+2\sqrt{2}+1\right|}{\sqrt{2}} \leq \left(\dfrac{b^2}{4}+1\right)$

        易知$\Delta = 1 - (2\sqrt{2}+1) < 0$，故不等式可以转化为$\dfrac{b^2}{4} - b+2\sqrt{2}+1 \leq \sqrt{2}\left(\dfrac{b^2}{4}+1\right)$，进一步化简有$\dfrac{(\sqrt{2}-1)b^2}{4}+b-\sqrt{2} -1 \geq 0$解得$b \in (-\infty, -(6+4\sqrt{2})] \cup [2, +\infty)$

        而动圆$C$的面积为$\left(\dfrac{b^2}{4}+1\right)^2\pi$，故其最小值为$4\pi$，当$b = 2$时取得最小值.

        \ps 求与已知直线平行的切线更方便些
    \end{jieda}
    \item 过抛物线 ${y}^{2} = {4x}$ 的焦点作一条倾斜角为 $\alpha$ ,长度不超过 8 的弦,弦所在的直线与圆 ${x}^{2} +$  ${y}^{2} = \dfrac{3}{4}$ 有公共点,则 $\alpha$ 的取值范围是\blkda[3]{$\left[\dfrac{\pi}{4}, \dfrac{\pi}{3}\right] \cup \left[\dfrac{2\pi}{3}, \dfrac{3\pi}{4}\right]$}.
    \begin{jieda}
        设弦所在直线为$y = k(x-1)$
        \begin{dingautolist}{172}
            \item 欲使弦与圆${x}^{2} +$  ${y}^{2} = \dfrac{3}{4}$ 有公共点，则原点到弦的距离应当小于等于半径，即$\dfrac{|k|}{\sqrt{1+k^2}} \leq \dfrac{\sqrt{3}}{2}$，故$k \in [-\sqrt{3}, \sqrt{3}]$

            \item $|AB| = \dfrac{2p}{sin^2\alpha} = \dfrac{4}{sin^2\alpha} \leq 8$，故$sin^2\alpha \geq \dfrac{1}{2}$，即$k^2 \geq 1$
        \end{dingautolist}
        综上，$k \in [-\sqrt{3}, -1] \cup [1, \sqrt{3}]$，因此$\alpha \in \left[\dfrac{\pi}{4}, \dfrac{\pi}{3}\right] \cup \left[\dfrac{2\pi}{3}, \dfrac{3\pi}{4}\right]$
    \end{jieda}
    
    \item 设直线 $l$ 与抛物线 ${y}^{2} = {4x}$ 相交于 $A\text{ 、 }B$ 两点,与圆 ${\left( x - 5\right) }^{2} + {y}^{2} = {r}^{2}\left( {r > 0}\right)$ 相切于点 $M$ , 且 $M$ 为线段 ${AB}$ 的中点,若这样的直线 $l$ 恰有 4 条,则 $r$ 的取值范围是\blkda{$(2, 4)$}.
    \begin{jieda}
        数形结合可知有两条切线为垂直于$x$轴的切线

        设斜率存在时$M(x_0, y_0)$，因此对应的切线$AB$为$(x-5)(x_0 - 5)+yy_0 = r^2$，设$A(x_1, y_1), B(x_2, y_2)$，则有$\left\{\begin{array}{l}
             y_1^2 = 4x_1 \\
             y_2^2 = 4x_2
        \end{array}\right.$，两者作差有$(y_1^2 - y_2^2) = 4(x_1-x_2)$即$(y_1+y_2)(y_1 - y_2) = 4(x_1-x_2)$，故$k_{AB} = \dfrac{4}{y_1+y_2} = \dfrac{2}{y_0}$，根据方程$(x-5)(x_0 - 5)+yy_0 = r^2$可知$k_{AB} = \dfrac{5-x_0}{y_0}$，因此$x_0 = 3$

        故当圆上存在横坐标为$3$的点时会出现另外两条切线，但根据图象可知圆上横坐标为$3$的点应在抛物线内部才可以，即$r \in (2, 4)$
    \end{jieda}
    
    \item 已知曲线 ${C}_{1} : y = {x}^{2} + 4$ 和 ${C}_{2} : y = {2x} - {x}^{2}$ ,直线 ${l}_{1}$ 与 ${C}_{1}\text{ 、 }{C}_{2}$ 分别相切于 $A\text{ 、 }B$ ,直线 ${l}_{2}$ (不同于 ${l}_{1}$ )与 ${C}_{1}$ 、 ${C}_{2}$ 分别相切于点 $C$ 、 $D$ ,则 ${AB}$ 与 ${CD}$ 交点的横坐标是\blkda{$\dfrac{1}{2}$}.
    \begin{jieda}
        $C_2: y = -(x-1)^2+1$，故$C_1, C_2$开口大小相同，存在对称中心$\left(\dfrac{1}{2}, \dfrac{5}{2}\right)$，公切线的交点应该是对称中心，此时两切线也各自关于该点中心对称.
    \end{jieda}
    
    \item 已知直线 $y = k\left( {x - m}\right)$ 与抛物线 ${y}^{2} = {2px}\left( {p > 0}\right)$ 交于 $A,B$ 两点, $O$ 为坐标原点,且 ${OA} \bot  {OB}$ ,又 ${OD} \bot  {AB}$ 于 $D$ ,若动点 $D$ 的坐标满足方程 ${x}^{2} + {y}^{2} - {4x} = 0$ ,则 $p = $ \blkda{2}.
    \begin{jieda}
        考虑特殊情况，当$AB$垂直于$x$轴时，要满足$OA \perp OB$，则应有$x_A = x_B = 2p$，此时$D(2p, 0)$，而${x}^{2} + {y}^{2} - {4x} = 0$与$x$轴的交点为$(0, 0)$和$(4,0)$，显然前者不符合要求，因此$p = 2$

        经检验$p = 2, m = 4$满足要求.
    \end{jieda}
    \item 已知点 $E$ 是抛物线 $C : {y}^{2} = {2px}\left( {p > 0}\right)$ 的对称轴与准线的交点,点 $F$ 为抛物线 $C$ 的焦点, 点 $P$ 在抛物线 $C$ 上,在 $\bigtriangleup {EFP}$ 中,若 $\sin \angle {EFP} = \mu  \cdot  \sin \angle {FEP}$ ,则 $\mu$ 的最大值为\blkda{$\sqrt{2}$}.
    \begin{jieda}
        $\mu = \dfrac{sin\angle{EFP}}{sin\angle{FEP}} = \dfrac{|EP|}{|FP|} = \dfrac{|EP|}{|PH|} = \dfrac{1}{sin\angle{PEH}}$，其中$H$是过$P$做准线的垂线的垂足

        因此当$P$为过$E$与抛物线相切的切点时，$\mu$最大，设$P(x_0, y_0)$，切线为$y = \dfrac{y_0}{2x_0}x+\dfrac{y_0}{2}$，令其过$E\left(-\dfrac{p}{2}, 0\right)$，可得$P\left(\dfrac{p}{2}, p\right)$，因此$\mu = \sqrt{2}$
    \end{jieda}
\end{enumerate}
\section*{11月6日}
\begin{enumerate}
    \item $y^2 = 4x$上两点$A, B$满足$\vv{OA}\cdot \vv{OB} = -3$，则$O$到$AB$的最大距离是\blkda{3}
    \begin{jieda}
        \begin{enumerate}
            \item 方法1

            设$A(x_1, y_1), B(x_2, y_2)$，则$x_1x_2+y_1y_2=-3$，即$\dfrac{(y_1y_2)^2}{16}+y_1y_2 = -3$，故$y_1y_2 = -4$或$-12$

            设直线 $AB$ 的方程为 $x = {m}y + {t}$ ,联立抛物线方程可得 ${y^2} - 4{m}y - 4{t} = 0$，则根据韦达定理$t = 1$或$3$，即直线恒过$(1, 0)$或$(3, 0)$，故$O$到$AB$的最大距离是$3$.

            \item 方法2

            抛物线的参数方程$\left\{\begin{array}{l}
                 x = 4t^2 \\
                 y = 4t
            \end{array}\right.$，则有$A(4t_1^2, 4t_1), B(4t_2^2, 4t_2)$，故$16t_1^2t_2^2+16t_1t_2 = -3$，故$t_1t_2 = -\dfrac{1}{4}$或$-\dfrac{3}{4}$

            设直线$AB: x = my+n$，将抛物线参数方程代入可得$4t^2 - 4tm-n = 0$，故根据韦达定理可知$-\dfrac{n}{4} = t_1t_2$，故直线$AB$恒过$(1, 0)$或$(3, 0)$，故$O$到$AB$的最大距离是$3$.
        \end{enumerate}
        
    \end{jieda}
    \item 自抛物线 ${y}^{2} = {4x}$ 上一点 $A\left( {1,2}\right)$ 引两弦 ${AM}\text{ 、 }{AN}$（$A$在$MN$上方） ,已知两弦的斜率之和为零, $\bigtriangleup  {AMN}$ 面积的最大值是\blkda{$\dfrac{32\sqrt{3}}{9}$}.
    \begin{jieda}
        $k_{MN} = -1$
        
        推导如下：设$M\left( {{x}_{1},{y}_{1}}\right) ,N\left( {{x}_{2},{y}_{2}}\right)$，则${k}_{1} = \dfrac{{y}_{1} - 2}{{x}_{1} - 1} = \dfrac{{y}_{1} - 2}{\dfrac{{y}_{1}^{2}}{4} - 1} = \dfrac{4}{{y}_{1} + 2}$ ,同理可得 ${k}_{2} = \dfrac{{y}_{2} - 2}{{x}_{2} - 1} = \dfrac{4}{{y}_{2} + 2}$. 
        由 ${k}_{1} + {k}_{2} = 0$ ,得 $\dfrac{4}{{y}_{1} + 2} + \dfrac{4}{{y}_{2} + 2} = 0$ ,即 ${y}_{1} + {y}_{2} + 4 = 0$. 
        设直线 $MN$ 的方程为 $x = {m}y + {t}$ ,联立抛物线方程可得 ${y^2} - 4{m}y - 4{t} = 0$. 
        则 $\Delta  > 0, \left\{\begin{array}{l}
             {y}_{1} + {y}_{2} = 4{m}  \\
             {y}_{1}{y}_{2} =  - 4{t} 
        \end{array}\right.$ 代入${y}_{1} + {y}_{2} + 4 = 0$可得 $4{m} + 4 = 0,{m} =  - 1$
        
        即$k_{MN} = -1$，$MN: x + y - {t} = 0$，$16+16t > 0, \left\{\begin{array}{l}
             {y}_{1} + {y}_{2} = -4  \\
             {y}_{1}{y}_{2} =  -4{t} 
        \end{array}\right.$，又因为$A$在$MN$上方，所以$t \in (-1, 3)$

        \begin{dingautolist}{172}
            \item 底乘高
            
            $A$到$MN$的距离为$d = \dfrac{|3-t|}{\sqrt{2}}$，故$S = \dfrac{1}{2} \times \dfrac{|3-t|}{\sqrt{2}} \times \sqrt{2}|y_1-y_2| = \dfrac{1}{2} \times |3-t| \times \sqrt{16+16t}$，$S^2 = 4(3-t)^2(1+t)$
            \item 利用邻边向量
            
            $\vv{AM} = (x_1-1, y_1-2), \vv{AN} = (x_2-1, y_2-2)$，$S = \dfrac{1}{2}|(x_1-1)(y_2-2)-(x_2-1)(y_2-2)| = \dfrac{1}{2}|(-y_1+t-1)(y_2-2)-(-y_2+t-1)(y_2-2)| = \dfrac{1}{2}|(3-t)(y_1-y_2)| = \dfrac{1}{2} \times |3-t| \times \sqrt{16+16t}$
        \end{dingautolist}
        
        \begin{dingautolist}{172}
            \item 基本不等式

            $4(3-t)^2(1+t) = 2(3-t)^2(2+2t) \leq 2\times\left(\dfrac{(3-t)+(3-t)+(2+2t)}{3}\right)^3 = \dfrac{1024}{27}$，当$t = \dfrac{1}{3}$时取等号，因此$S \leq \dfrac{32\sqrt{3}}{9}$
        
            \item 求导

            设$f(x) = 4(x-3)^2(1+x), x > -1$，$f'(x) = 8(x-3)(1+x) + 4(x-3)^2 = 4(x-3)(3x-1)$，故$f(x)$在$x = \dfrac{1}{3}$处取得最大值，即$S^2 \leq \dfrac{1024}{27}$，故$S \leq \dfrac{32\sqrt{3}}{9}$
        \end{dingautolist}
    \end{jieda}
    \item 自抛物线 ${y}^{2} = {4x}$ 上一点 $A\left( {1,2}\right)$ 引两弦 ${AM}\text{ 、 }{AN}$ ,已知两弦的斜率之积为$-4$, 则$A$到$MN$的距离的最大值为\blkda{$\sqrt{17}$}.
    \begin{jieda}
        $MN$过定点$(2, -2)$

        推导如下：设$M\left( {{x}_{1},{y}_{1}}\right) ,N\left( {{x}_{2},{y}_{2}}\right)$，则${k}_{1} = \dfrac{{y}_{1} - 2}{{x}_{1} - 1} = \dfrac{{y}_{1} - 2}{\dfrac{{y}_{1}^{2}}{4} - 1} = \dfrac{4}{{y}_{1} + 2}$ ,同理可得 ${k}_{2} = \dfrac{{y}_{2} - 2}{{x}_{2} - 1} = \dfrac{4}{{y}_{2} + 2}$. 
        由 ${k}_{1} \cdot {k}_{2} = -4$ ,得 $\dfrac{4}{{y}_{1} + 2} \cdot \dfrac{4}{{y}_{2} + 2} = -4$ ,即 $y_1y_2+2(y_1+y_2)+4 = -4$. 
        设直线 $MN$ 的方程为 $x = {m}y + {t}$ ,联立抛物线方程可得 ${y^2} - 4{m}y - 4{t} = 0$. 
        则 $\Delta  > 0, \left\{\begin{array}{l}
             {y}_{1} + {y}_{2} = 4{m}  \\
             {y}_{1}{y}_{2} =  - 4{t} 
        \end{array}\right.$ 代入$y_1y_2+2(y_1+y_2)+4 = -4$可得 $-4t+2(4m)+4 = -4$，即$t-2m = 2$，因此$MN$恒过$(2, -2)$

        故$A$到$MN$的距离的最大值为$\sqrt{17}$
    \end{jieda}
    
    \item 过原点 $O$ 作两条相互垂直的直线分别与椭圆 $P : \dfrac{{x}^{2}}{2} + {y}^{2} = 1$ 交于 $A\text{ 、 }C$ 与 $B\text{ 、 }D$ ,则四边形 ${ABCD}$ 面积的最小值为\dotfill{(\kh{A})}
    \xx{$\dfrac{8}{3}$}{$4\sqrt{2}$}{$2\sqrt{2}$}{$\dfrac{4}{3}$}
    \begin{jieda}
        \begin{dingautolist}{172}
            \item 当两直线分别与坐标轴平行时，$S = 2\sqrt{2}$
            \item 当两直线斜率均存在时，设$AC$为$y = kx$，与椭圆方程联立得$\left\{\begin{array}{l}
                 y = kx  \\
                 \dfrac{{x}^{2}}{2} + {y}^{2} = 1 
            \end{array}\right.$，消去$y$得到$(1+2k^2)x^2 = 2$，因此$|x_1 - x_2| = 2\sqrt{\dfrac{2}{1+2k^2}}$，则$|AC| = 2\sqrt{\dfrac{2+2k^2}{1+2k^2}}$，同理可得$|BD| = 2\sqrt{\dfrac{2+2k^2}{2+k^2}}$，故$S = \dfrac{1}{2}|AC| \cdot |BD| = 4\sqrt{\dfrac{1+k^2}{1+2k^2} \cdot \dfrac{1+k^2}{2+k^2}}$

            令$t = k^2+1, t \in [1, +\infty)$，则$S = 4\sqrt{\dfrac{t^2}{2t^2+t-1}} = \dfrac{4}{\sqrt{2+\dfrac{1}{t}-\dfrac{1}{t^2}}} \geq \dfrac{8}{3}$
        \end{dingautolist}
    \end{jieda}
    \item 设 $a\text{ 、 }b$ 是两个实数,且满足 $0 \leq  a < b$ ,直线 $l : y = {kx} + m$ 和圆 ${x}^{2} + {y}^{2} = 1$ 交于两点 $A\text{ 、 }B$ ,若存在正数 $m$ ,对于任意的 $k \in \left\lbrack  {a,b}\right\rbrack$ ,均使得 $\bigtriangleup {OAB}$ 的面积不小于 $\dfrac{\sqrt{3}}{4}$ ,则 $b - {2a}$ 的最大值为\blkda{$\sqrt{2}$}.
    \begin{jieda}
        设 $O$ 到直线 $l$ 的距离为 $d$ ,则${S}_{\bigtriangleup {AOB}} = \dfrac{1}{2} \times  2\sqrt{1 - {d}^{2}} \times  d \geq  \dfrac{\sqrt{3}}{4}$，解得 $\dfrac{1}{2} \leq  d \leq  \dfrac{\sqrt{3}}{2}$，即 $\dfrac{1}{2} \leq  \dfrac{m}{\sqrt{{k}^{2} + 1}} \leq  \dfrac{\sqrt{3}}{2}$，所以$\dfrac{1}{2}\sqrt{{k}^{2} + 1} \leq  m \leq  \dfrac{\sqrt{3}}{2}\sqrt{{k}^{2} + 1}$，若存在正数 $m$，对于任意的 $k \in \left\lbrack  {a,b}\right\rbrack$均满足要求，则应有${\left( \dfrac{1}{2}\sqrt{{k}^{2} + 1}\right) }_{\text{max }} \leq m \leq {\left( \dfrac{\sqrt{3}}{2}\sqrt{{k}^{2} + 1}\right) }_{\text{min }}$

        因为$\left\{\begin{array}{l}
             {\left( \dfrac{1}{2}\sqrt{{k}^{2} + 1}\right) }_{\text{max }} = \dfrac{1}{2}\sqrt{{b}^{2} + 1}  \\
             {\left( \dfrac{\sqrt{3}}{2}\sqrt{{k}^{2} + 1}\right) }_{\text{min }} = \dfrac{\sqrt{3}}{2}\sqrt{{a}^{2} + 1} 
        \end{array}\right.$，所以应有$\dfrac{1}{2}\sqrt{{b}^{2} + 1} \leq  \dfrac{\sqrt{3}}{2}\sqrt{{a}^{2} + 1}$ ,化简得$\dfrac{{b}^{2}}{2} - \dfrac{3{a}^{2}}{2} \leq  1$ ,且 $0 \leq  a < b$
        \begin{center}
            \includegraphics[width=0.4\linewidth]{bodyxb//src/46-12.png}
        \end{center}
        设 $b - {2a} = t$ ,即 $b = {2a} + t$，在平面直角坐标系 $a - O - b$ 中,直线$b = {2a} + t$ 和由 $\dfrac{{b}^{2}}{2} - \dfrac{3{a}^{2}}{2} \leq  1$，且 $0 \leq  a < b$ 表示的平面区域有公共点
        
        当直线 $b = {2a} + t$ 经过点 $A(0$ , $\sqrt{2})$ 时, $t$最大, ${t}_{max} = \sqrt{2}$ .
    \end{jieda}
\end{enumerate}
\section*{11月13日}
\begin{enumerate}

    \item 已知 $\sqrt{{x}^{2} + {y}^{2}} + \sqrt{{\left( x - 8\right) }^{2} + {\left( y - 6\right) }^{2}} = {20}$ ,则 $\left| {{3x} - {4y} - {100}}\right|$ 的最大值为\blkda[3]{$100+25\sqrt{3}$}.
    \begin{jieda}
        根据椭圆定义可知$\sqrt{{x}^{2} + {y}^{2}} + \sqrt{{\left( x - 8\right) }^{2} + {\left( y - 6\right) }^{2}} = {20}$对应着以$(0, 0), (8, 6)$为焦点，长轴长为$20$的椭圆

        $\left| {{3x} - {4y} - {100}}\right|$即点$(x, y)$到直线$3x - {4y} - 100 = 0$距离的5倍

        易知椭圆$\sqrt{{x}^{2} + {y}^{2}} + \sqrt{{\left( x - 8\right) }^{2} + {\left( y - 6\right) }^{2}} = {20}$的长轴与直线$3x - {4y} - 100 = 0$平行，故其最大值在椭圆短轴顶点处取到，为椭圆中心到直线距离加短半轴长的和的5倍，为$100+25\sqrt{3}$
    \end{jieda}
    \item 已知椭圆方程为 $\dfrac{{x}^{2}}{2} + {y}^{2} = 1$ ,斜率为 1 的直线 $l$ 与椭圆交于 $A,B$ 两个不同的点, ${AB}$ 中点为 $Q$ ,点 $P\left( {2,0}\right)$ ,则 ${\left| PQ\right| }^{2}$ 的取值范围是\blkda[4]{$\left(\dfrac{17-8\sqrt{3}}{3}, \dfrac{17+8\sqrt{3}}{3}\right)$}.
    \begin{jieda}
        设$A(x_1, y_2), B(x_2, y_2)$，则$\left\{\begin{array}{l}
             \dfrac{{x_1}^{2}}{2} + {y_1}^{2} = 1 \\
             \dfrac{{x_2}^{2}}{2} + {y_2}^{2} = 1 
        \end{array}\right.$，两式做差可得$\dfrac{x_1^2-x_2^2}{2} + (y_1^2-y_2^2) = 0$，即$\dfrac{(x_1+x_2)(x_1-x_2)}{2} + (y_1+y_2)(y_1-y_2) = 0$，因为$\dfrac{y_1-y_2}{x_1-x_2} = 1$，所以$y_1+y_2 = -\dfrac{x_1+x_2}{2}$，即$Q$在$y = -\dfrac{x}{2}, x \in \left(-\dfrac{2\sqrt{3}}{3}, \dfrac{2\sqrt{3}}{3}\right)$上，设$Q(x, y)$，则$|PQ|^2 = (x-2)^2+y^2 = \dfrac{5}{4}x^2-4x+4$
        易知函数$f(x) = \dfrac{5}{4}x^2-4x+4$在$\left(-\dfrac{2\sqrt{3}}{3}, \dfrac{2\sqrt{3}}{3}\right)$上严格减，因此$f(x) \in \left(\dfrac{17-8\sqrt{3}}{3}, \dfrac{17+8\sqrt{3}}{3}\right)$
    \end{jieda}
    \item 直线 $\sqrt{3}{ax} + {by} = 1\left( {a,b \in  \mathbf{R}}\right)$ 与圆 $O : {x}^{2} + {y}^{2} = 2$ 相交于 $A,B$ 两点,且 $\bigtriangleup {AOB}$ 是直角三角形 $(O$ 是坐标原点),则点 $P\left( {a,b}\right)$ 与点 $Q\left( {0,1}\right)$ 之间距离的最大值是\blkda{2}.
    \begin{jieda}
        因为$\bigtriangleup {AOB}$ 是直角三角形，$A,B$在圆上，因此只有$\angle{AOB}$可能为直角，故原点到直线 $\sqrt{3}{ax} + {by} = 1\left( {a,b \in  \mathbf{R}}\right)$的距离为$1$，即$\dfrac{1}{\sqrt{3a^2+b^2}} = 1$，所以$3a^2+b^2 = 1$，即$P$在椭圆$3x^2 + y^2 = 1$上，因此$PQ$的最大值为2
    \end{jieda}
    \item 如图,已知 $F$ 是椭圆 $\dfrac{{x}^{2}}{4} + \dfrac{{y}^{2}}{3} = 1$ 的左焦点, $A$ 为椭圆的下顶点,点 $P$ 是椭圆上任意一点,以 ${PF}$ 为直径作圆 $N$ ,射线 ${ON}$ 与圆 $N$ 交于点 $Q$ ,则 $\left| {AQ}\right|$ 的取值范围为\blkda[3]{$\left\lbrack  {2 - \sqrt{3},2 + \sqrt{3}}\right\rbrack$}.
    
    \begin{center}
    \includegraphics[width=0.3\textwidth]{bodyxb/src/bo_d1ig044601uc738hj6v0_93_1055_1304_435_325_0.jpg}
    \end{center}
    \begin{jieda}
        设 ${F}_{1}$ 为椭圆的右焦点，因为 $O$ , $N$ 分别为 $F{F}_{1}$, ${PF}$ 的中点，所以 $\left| {ON}\right|  = \dfrac{1}{2}\left| {P{F}_{1}}\right|$，$\left|{NQ}\right|  =  \dfrac{1}{2}\left| {PF}\right|$，因此$\left| {OQ}\right|  = \left| {ON}\right|  + \left| {NQ}\right|  = \dfrac{1}{2}\left( {\left| {PF}\right|  + \left| {P{F}_{1}}\right| }\right)  = 2$
        
        所以 $Q$ 的轨迹是以 $O$ 为圆心,以2为半径的圆, $A(0, - \sqrt{3})$ 在圆内，所以数形结合可知$\left| {AQ}\right|$ 的取值范围为 $\left\lbrack  {2 - \sqrt{3},2 + \sqrt{3}}\right\rbrack$ .
    \end{jieda}
    \item 若对圆 ${\left( x - 3\right) }^{2} + {\left( y - 2\right) }^{2} = 1$ 上任意一点 $P\left( {x,y}\right) ,\left| {{3x} - {4y} + a}\right|  + \left| {9 - {3x} + {4y}}\right|$ 的取值与 $x,y$ 无关,则实数 $a$ 的取值范围是\blkda{$[4, +\infty)$}.
    \begin{jieda}
        设$k = {3x} - {4y}$，则$\left| {{3x} - {4y} + a}\right|  + \left| {9 - {3x} + {4y}}\right| = |k+a|+|9-k|$，易知其在$[-a, 9]$或$[9, -a]$范围内值不变

        因为$P$在圆上，因此数形结合可知$k \in [-4, 6]$（相切时取得最值），故$[-4, 6] \subseteq [-a, 9]$，因此$a \geq 4$
    \end{jieda}
    \item 已知过原点的直线 $l$ 交椭圆 $\dfrac{{x}^{2}}{9} + \dfrac{{y}^{2}}{4} = 1$ 于 $P,Q$ 两点,椭圆右顶点为 $A,\bigtriangleup {APQ}$ 重心为 $G$ , $\left| {GA}\right|  + \left| {GP}\right|  + \left| {GQ}\right|$ 的取值范围是\blkda[3]{$[2\sqrt{5}+2, 8)$}.
    \begin{jieda}
        $A(3, 0)$，根据重心性质可知$G(1, 0)$，设$P(x_0, y_0)$，则$Q(-x_0, y_0)$，故$|GP|+|GQ| = \sqrt{(x_0-1)^2+y_0^2}+\sqrt{(x_0+1)^2+y_0^2}$

        当$\sqrt{(x_0-1)^2+y_0^2}+\sqrt{(x_0+1)^2+y_0^2}$等于大于2的常值时，$(x_0, y_0)$在以$(1, 0), (-1, 0)$为焦点的椭圆$\dfrac{x^2}{k}+\dfrac{y^2}{k-1} = 1$上，故数形结合可知$k \in [5, 9)$（$k = 9$时$Q$与$A$重合），即$\sqrt{(x_0-1)^2+y_0^2}+\sqrt{(x_0+1)^2+y_0^2} \in [2\sqrt{5}, 6)$

        $\left| {GA}\right|  + \left| {GP}\right|  + \left| {GQ}\right| \in [2\sqrt{5}+2, 8)$
    \end{jieda}
\end{enumerate}
\section*{12月4日}
\begin{enumerate}
    \item 如图,正方体的棱长是 $a$ ,点 $E\text{ 、 }F$ 分别是两条棱的中点. 求证:四边形 ${BDEF}$ (图中阴影部分)是一个梯形
    
    \begin{center}
    \includegraphics[width=0.2\textwidth]{bodyxb/src/bo_d4lnl5ref24c739do320_11_241_208_336_326_0.jpg}
    \end{center}
    \begin{jieda}
        如图所示，连接 ${B}_{1}{D}_{1}$
        
        \begin{center}
        \includegraphics[width=0.2\textwidth]{bodyxb/src/bo_d4lnl5ref24c739do320_11_243_1097_335_324_0.jpg}
        \end{center}
        
        
        因为点 $E$ 、 $F$ 分别是 ${A}_{1}{D}_{1}$ 、 ${A}_{1}{B}_{1}$ 的中点，所以 ${EF}//{B}_{1}{D}_{1}$ ， ${EF} = \dfrac{1}{2}{B}_{1}{D}_{1}$ ，
        
        又因为 ${BD}//{B}_{1}{D}_{1},{BD} = {B}_{1}{D}_{1}$ ,所以 ${EF}//{BD},{EF} = \dfrac{1}{2}{BD}$ ,
        
        所以四边形 ${BDEF}$ 是一个梯形.
    \end{jieda}
    \item 如图所示，在空间四边形 ${ABCD}$ 中，点 $E$ ， $H$ 分别是边 ${AB}$ ， ${AD}$ 的中点，点 $F$ ， $G$ 分别是边 ${BC},{CD}$ 上的点，且 $\dfrac{CF}{CB} = \dfrac{CG}{CD} = \dfrac{2}{3}$ ，有以下结论正确的是\blkda{\ding{173}\ding{175}}
    
    \begin{center}
    \includegraphics[width=0.2\textwidth]{bodyxb/src/bo_d4lnl5ref24c739do320_2_243_995_278_307_0.jpg}
    \end{center}

    \begin{dingautolist}{172}
        \item $EF$与$GH$平行
        \item $EF$与$GH$共面
        \item $EF$与$GH$的交点 $M$ 可能在直线 ${AC}$ 上，也可能不在直线 ${AC}$ 上
        \item ${EF}$ 与 ${GH}$ 的交点 $M$ 一定在直线 ${AC}$ 上.
    \end{dingautolist}

    \begin{jieda}
        如图所示. 连接 ${EH},{FG}$ ,
        
        \begin{center}
        \includegraphics[width=0.2\textwidth]{bodyxb/src/bo_d4lnl5ref24c739do320_3_239_193_280_308_0.jpg}
        \end{center}
        
        
        依题意,可得 ${EH}//{BD},{FG}//{BD}$，所以 ${EH}//{FG}$，所以 $E,F,G,H$ 共面
        
        因为 ${HE} = \dfrac{1}{2}{BD},{FG} = \dfrac{2}{3}{BD}$，所以四边形 ${EFGH}$ 是梯形, ${EF}$ 与 ${GH}$ 必相交，设EF与 ${GH}$ 的交点为 $M$
        
        因为点 $M$ 在 ${EF}$ 上，故点 $M$ 在平面 ${ACB}$ 上；同理，点 $M$ 在平面 ${ACD}$ 上,
        
        所以点 $M$ 在平面 ${ACB}$ 与平面 ${ACD}$ 的交线上， 又 ${AC}$ 是这两个平面的交线，所以点 $M$ 一定在直线 ${AC}$ 上
    \end{jieda}
    
    \item 如图所示,在四棱锥 $P - {ABCD}$ 中, ${BC}//$ 平面 ${PAD},{BC} = \dfrac{1}{2}{AD},E$ 是 ${PD}$ 的中点.
    
    \begin{center}
    \includegraphics[width=0.3\textwidth]{bodyxb//src/bo_d4lsgh3ef24c739dpshg_2_240_1723_425_416_0.jpg}
    \end{center}
    \begin{enumerate}
        \item 求证: ${BC}//{AD}$
        \item 求证: ${CE}//$ 平面 ${PAB}$
        \item *若 $M$ 是线段 ${CE}$ 上一动点，则线段 ${AD}$ 上是否存在点 $N$ ，使 ${MN}//$ 平面 ${PAB}$ ？说明理由
    \end{enumerate}
    \begin{jieda}
        \begin{enumerate}
            \item 在四棱锥 $P - {ABCD}$ 中， ${BC}//$ 平面 ${PAD}$ ， ${BC} \subset$ 平面 ${ABCD}$ ， ${AD} \subset$ 平面 ${PAD}$，平面 ${ABCD} \cap$ 平面 ${PAD} = {AD}$ ,所以 ${BC}//{AD}$
            
            \item 如下图，取 $F$ 为 ${AP}$ 中点，连接 ${EF},{BF}$ ，由 $E$ 是 ${PD}$ 的中点，

            所以 ${EF}//{AD}$ 且 ${EF} = \dfrac{1}{2}{AD}$ ,由上一问知 ${BC}//{AD}$ ,又 ${BC} = \dfrac{1}{2}{AD}$ ,
            
            所以 ${EF}//{BC}$ 且 ${EF} = {BC}$ ,所以四边形 ${BCEF}$ 为平行四边形,故 ${CE}//{BF}$ ,
            
            而 ${CE} \subset$ 平面 ${PAB},{BF} \not \subset $ 平面 ${PAB}$ ,则 ${CE}//$ 平面 ${PAB}$ .
            
            \begin{center}
            \includegraphics[width=0.3\textwidth]{bodyxb//src/bo_d4lsgh3ef24c739dpshg_3_243_1327_434_427_0.jpg}
            \end{center}

            \item 取 ${AD}$ 中点 $N$ ，连接 ${CN}$ ， ${EN}$ ，
            
            因为 $E,N$ 分别为 ${PD},{AD}$ 的中点,所以 ${EN}//{PA}$ ,
            
            因为 ${EN} \not \subset$ 平面 ${PAB},{PA} \subset$ 平面 ${PAB}$ ,所以 ${EN}//$ 平面 ${PAB}$ ,
            
            线段 ${AD}$ 存在点 $N$ ，使得 ${MN}//$ 平面 ${PAB}$ ，理由如下:
            
            由上一问知: ${CE}//$ 平面 ${PAB}$ ，又 ${CE} \cap  {EN} = E$ ， ${CE} \subset$ 平面 ${CEN}$ ， ${EN} \subset$ 平面 ${CEN}$ ， 所以平面 ${CEN}//$ 平面 ${PAB}$ ,又 $M$ 是 ${CE}$ 上的动点, ${MN} \subset$ 平面 ${CEN}$ ,
            
            所以 ${MN}//$ 平面 ${PAB}$ ,所以线段 ${AD}$ 存在点 $N$ ,使得 ${MN}//$ 平面 ${PAB}$ .
        \end{enumerate}
    \end{jieda}
    \item 求证：已知三个平面两两相交，且有三条交线. 如果这三条交线中有两条交线平行，那么第三条交线也与另两条直线平行.
    \begin{jieda}
        设$\alpha \cap \beta = a, \beta \cap \gamma = b, \alpha \cap \gamma = c$

        $a // b$则根据线面平行的判定定理可知$a // \beta$，故根据线面平行性质定理可知$a // c$，根据公理4可知$b //c$.
    \end{jieda}     
    \item 如图，在四棱锥 $P - {ABCD}$ 中，底面 ${ABCD}$ 是边长为 $a$ 的正方形， $O$ 为正方形 ${ABCD}$ 的中心， ${PO} \bot$ 平面 ${ABCD}$ . 若点 $M$ 在棱 ${PC}$ 上且不与 $P$ 、 $C$ 重合，平面 ${ADM}$ 交棱 ${PB}$ 于点 $N$ ，求证: ${MN}//{AD}$ .
    
    \begin{center}
    \includegraphics[width=0.3\textwidth]{bodyxb/src/bo_d4lnl5ref24c739do320_12_241_500_373_340_0.jpg}
    \end{center}
    \begin{jieda}
        因为${AD}//{BC}$ 且 ${AD} \subset$ 平面 ${ADM}$ ， ${BC}$ 不在平面 ${ADM}$ 上， 所以${BC}//$ 平面 ${ADM}$ ，

        又 因为${BC} \subset$ 平面 ${PBC}$ ,平面 ${ADM} \cap$ 平面 ${PBC} = {MN}$，因此$ {BC}//{MN}$ ,且 ${BC}//{AD}$，所以${MN}//{AD}$ .
    \end{jieda}
    \item 如图所示，已知三棱柱 ${ABC} - {A}^{\prime }{B}^{\prime }{C}^{\prime }$ 中，$D$ 是 $BC$ 的中点，$D'$是 $B{C}^{\prime }$ 的中点，设平面 ${A}^{\prime }{D}^{\prime }B \cap$ 平面 ${{ABC}} = {a}$ ,平面 ${{AD}}{{C}}^{\prime } \cap$ 平面 ${{A}}^{\prime }{{B}}^{\prime }{{C}}^{\prime } = {b}$ ,判断直线 ${a},{b}$ 的位置关系,并证明.
    
    \begin{center}
    \includegraphics[width=0.2\textwidth]{bodyxb/src/bo_d4lnl5ref24c739do320_4_265_1082_306_354_0.jpg}
    \end{center}
    \begin{jieda}
        连接$DD'$，因为$D$ 是 $BC$ 的中点，$D'$是 $B{C}^{\prime }$ 的中点，所以$DD' // BB'$，且$DD' = BB'$，故$DD' // AA'$，且$DD' = AA'$，故$A'D' // AD$，且$A'D' = AD$，故$AD // $平面$A'B'C'$，$A'D' // $平面$ABC$
    
        因为平面 ${{A}}^{\prime }{{D}}^{\prime }{B} \cap$ 平面 ${{ABC}} = {a}$，$A'D' // $平面$ABC$，所以${{A}}^{\prime }{{D}}^{\prime }//{a}$，同理可得 ${{AD}}//{b}$，根据平行的传递性可知 ${a}//{b}$ .
    \end{jieda}
    \item 如图，在四棱锥 $P - {ABCD}$ 中，底面 ${ABCD}$ 是正方形. ${PD} \bot$ 平面 ${ABCD}$ ， ${PD} = {AD} = 2$ ，点 $E$ 是棱 ${PC}$ 的中点,平面 ${ABE}$ 与棱 ${PD}$ 交于点 $F$，求四边形$ABEF$的面积.
    
    \begin{center}
    \includegraphics[width=0.2\textwidth]{bodyxb/src/bo_d4lnl5ref24c739do320_13_241_198_342_344_0.jpg}
    \end{center}
    \begin{jieda}
        ${CD}//{AB},{AB} \subset$ 平面 ${ABE},{CD} \subset$ 平面 ${ABE}$，所以${CD}//$ 平面 ${ABE}$
        
        又 ${CD} \subset$ 平面 ${PCD}$ ,平面 ${PCD} \cap$ 平面 ${ABE} = {EF}$，故${CD}//{EF}$，故$EF$为$\triangle PAC$的中位线，所以$EF = 1$
        
        由于$CD//AB$，根据平行的传递性可知${AB}//{EF}$，故$ABEF$为梯形

        根据三垂线定理可知$AB \bot AF$，故$ABEF$为直角梯形，$EF = 1, AB = 2, AF = \sqrt{5}$，故面积为$\dfrac{3\sqrt{5}}{2}$.
    \end{jieda}    
    \item 在五面体 ${ABCDEF}$ 中, ${CD} \bot$ 平面 ${ADE},{EF} \bot$ 平面 ${ADE}$，求证: ${EF}//{AB}$
    
    \begin{center}
    \includegraphics[width=0.35\textwidth]{bodyxb/src/bo_d4lnl5ref24c739do320_9_243_1790_361_191_0.jpg}
    \end{center}
    \begin{jieda}
        因为 ${CD} \bot$ 平面 ${ADE}$ ， ${EF} \bot$ 平面 ${ADE}$ ，则 ${CD}//{EF}$ ，
        
        因为 ${EF} \not \subset$ 平面 ${ABCD},{CD} \subset$ 平面 ${ABCD}$ ,所以, ${EF}//$ 平面 ${ABCD}$ ,
        
        因为 ${EF} \subset$ 平面 ${ABFE}$ ,平面 ${ABFE} \cap$ 平面 ${ABCD} = {AB}$ ,所以, ${EF}//{AB}$ .
    \end{jieda}    
    
    \item 如图所示,在正方体 ${ABCD} - {A}_{1}{B}_{1}{C}_{1}{D}_{1}$ 中, ${EF}$ 与 ${AC},{A}_{1}D$ 都垂直相交,垂足分别是点 $F$ 、点 $E$ .

    \begin{center}
    \includegraphics[width=0.35\textwidth]{bodyxb//src/bo_d4lsgh3ef24c739dpshg_4_242_411_406_355_0.jpg}
    \end{center}
    \begin{enumerate}
        \item 求证: ${EF}\bot$ 平面 $A{B}_{1}C$
        \item 求证: ${EF}//{B{D}_{1}}$
    \end{enumerate}
    \begin{jieda}
        \begin{enumerate}
            \item 在正方体 ${ABCD} - {A}_{1}{B}_{1}{C}_{1}{D}_{1}$ 中， ${CD}//{A}_{1}{B}_{1}$ 且 ${CD} = {A}_{1}{B}_{1}$ ，
            
            所以四边形 ${A}_{1}{B}_{1}{CD}$ 是平行四边形,所以 ${A}_{1}D//{B}_{1}C$ ,
            
            因为 ${EF} \bot  {A}_{1}D$ ,所以 ${EF} \bot  {B}_{1}C$ ,
            
            因为 ${EF} \bot  {AC},{AC} \cap  {B}_{1}C = C,{B}_{1}C \subset$ 平面 $A{B}_{1}C,{AC} \subset$ 平面 $A{B}_{1}C$ .
            
            所以 ${EF} \bot$ 平面 $A{B}_{1}C$ .
            
            \begin{center}
            \includegraphics[width=0.3\textwidth]{bodyxb//src/bo_d4lsgh3ef24c739dpshg_4_244_1688_403_357_0.jpg}
            \end{center}
            \item 因为 $D{D}_{1} \bot$ 平面 ${ABCD}$ ， ${AC} \subset$ 平面 ${ABCD}$，所以 $D{D}_{1} \bot  {AC}$ ,
            
            又因为 ${BD} \bot  {AC},D{D}_{1} \cap  {BD} = D,{BD} \subset$ 平面 ${BD}{D}_{1}{B}_{1},D{D}_{1} \subset$ 平面 ${BD}{D}_{1}{B}_{1}$ , 所以 ${AC} \bot$ 平面 ${BD}{D}_{1}{B}_{1}$ ,
            
            因为 $B{D}_{1} \subset$ 平面 ${BD}{D}_{1}{B}_{1}$ ,所以 ${AC} \bot  B{D}_{1}$ ,
            
            同理可证 $B{D}_{1} \bot  {B}_{1}C$ ,
            
            又 ${AC} \cap  {B}_{1}C = C,{B}_{1}C \subset$ 平面 $A{B}_{1}C,{AC} \subset$ 平面 $A{B}_{1}C$ .
            
            所以 $B{D}_{1} \bot$ 平面 $A{B}_{1}C$ .
            
            又 ${EF} \bot$ 平面 $A{B}_{1}C$ ,所以 ${EF}//B{D}_{1}$ .
        \end{enumerate}
    \end{jieda}
\end{enumerate}

\section*{12月11日}
\begin{enumerate}
    \item 在 ${{Rt}} \bigtriangleup  {ACB}$ 中,已知 $\angle {ACB} = {90}^{ \circ  },{AC} = {BC} = 1,{PA}\bot$ 平面 ${ABC}$ ,且 ${PA} = \sqrt{2}$ (如图),则$cos\angle{PBA} = $\blkda{$\dfrac{\sqrt{2}}{2}$}，$cos\angle{PBC} = $\blkda{$\dfrac{1}{2}$}，$cos\angle{ABC} = $\blkda{$\dfrac{\sqrt{2}}{2}$}；
    设$Q$为面$ABC$内动点，则异面直线$CQ$与$PB$的夹角取值范围是\blkda{$\left[\dfrac{\pi}{4}, \dfrac{\pi}{2}\right]$}
    \begin{center}
    \includegraphics[width=0.3\textwidth]{bodyxb/src/bo_d20aucref24c73anc48g_148_454_647_255_218_0.jpg}
    \end{center}
    \begin{jieda}
        $PB$在面内射影为$AB$，$cos\angle{PBC} = cos\angle{ABC} \cdot cos\angle{PBA}$
    \end{jieda}
    \item ${PA},{PB},{PC}$ 是从点 $P$ 出发的不共面的三条射线,每两条射线的夹角均为 ${60}^{ \circ  }$ ,那么直线 ${PC}$ 与平面 ${PAB}$ 所成角的余弦值是\blkda{$\dfrac{\sqrt{3}}{3}$}
    \begin{jieda}
        如图,设直线 ${PC}$ 在平面 ${PAB}$ 的射影为 ${PD}$ ,
        
        \begin{center}
        \includegraphics[width=0.2\textwidth]{bodyxb/src/bo_d4ojskjef24c73bc0dj0_1_243_274_288_325_0.jpg}
        \end{center}

        根据三余弦定理可知$cos\angle{APC} = cos\angle{APD} \cdot cos\angle{CPD}$

        易知$\angle{APC} = \ang{60}$，$\angle{APD} = \ang{30}$，故$cos\angle{CPD} = \dfrac{\sqrt{3}}{3}$.
    \end{jieda}
    
    \item 已知三棱柱 ${ABC} - {A}_{1}{B}_{1}{C}_{1}$ 的侧棱与底面边长都相等， ${A}_{1}$ 在底面 ${ABC}$ 上的射影为底面的中心，则 ${AB_1}$ 与 底面$ABC$ 所成的角的正弦值为\blkda{$\dfrac{\sqrt{2}}{3}$}
    \begin{jieda}
        设底面的中心为 $O$，由题意可知 ${A}_{1}O \bot$ 平面 ${ABC}$
        
        设三棱柱 ${ABC} - {A}_{1}{B}_{1}{C}_{1}$ 的侧棱与底面边长为 1，根据几何关系可知$AO = \dfrac{\sqrt{3}}{3}$，故$A_1O = \dfrac{\sqrt{6}}{3}$，故$cos\angle{A_1AO} = \dfrac{\sqrt{3}}{3}$，又因为$\angle{BAO} = \dfrac{\pi}{6}$，故根据三余弦定理可知$cos\angle{A_1AB} = \dfrac{1}{2}$.

        因此$\angle{A_1AB} = \ang{60}$，故$AB_1 = 2\sqrt{3}$，所以所求角的正弦值为$\dfrac{B_1O_1}{AB_1} = \dfrac{A_1O_1}{AB_1} = \dfrac{\sqrt{2}}{3}$
    \end{jieda}
    
    \item 如图所示，在侧棱垂直于底面的三棱柱 ${ABC} - {A}_{1}{B}_{1}{C}_{1}$ 中， $P$ 是棱 ${BC}$ 上的动点，记直线 ${A}_{1}P$ 与平面 ${ABC}$ 所成的角为 ${\theta }_{1}$ ,与直线 ${BC}$ 所成的角为 ${\theta }_{2}$ ,则 ${\theta }_{1},{\theta }_{2}$ 的大小关系是
    
    \begin{center}
    \includegraphics[width=0.3\textwidth]{bodyxb/src/bo_d4ojskjef24c73bc0dj0_0_242_626_445_446_0.jpg}
    \end{center}

    \xx{${\theta }_{1} = {\theta }_{2}$}{${\theta }_{1} > {\theta }_{2}$}{${\theta }_{1} < {\theta }_{2}$}{不能确定}

    \begin{jieda}
        根据线面角是该直线与对应平面内的任意直线所成角中最小的角，所以有 ${\theta }_{1} < {\theta }_{2}$ ,故选 C.
    \end{jieda}

    
    \item 已知等腰$Rt \bigtriangleup {ABC}$ 内接于圆 $O$，连接$OC$，交点圆 $O$于$C'$， $M$ 是下半圆弧上的动点 (如图所示),现将上半圆面沿 ${AB}$ 折起，使得$\angle{COC'} = \dfrac{\pi }{4}$ . 则直线 ${AC}$ 与直线 ${OM}$ 所成角的最小值是\blkda{$\dfrac{\pi }{6}$}
    
    \begin{center}
    \includegraphics[width=0.25\textwidth]{bodyxb/src/bo_d4onr53ef24c73bc3ut0_0_246_691_285_267_0.jpg}
    \end{center}
    \begin{jieda}
        ${AC}$ 与直线 ${OM}$ 所成角的最小值是$AC$与平面$ABM$的线面所成角

        做$CH \bot $面$AOM$，则$CH = \dfrac{r}{\sqrt{2}}$，$AC = \sqrt{2}r$，故线面所成角为$\dfrac{\pi}{6}$.
    \end{jieda}
    
    \item 在立方体$ABCD-A_1B_1C_1D_1$中，$P$为底面$ABCD$内的动点，异面直线$BD_1$和$DP$夹角的最小值为\blkda[2]{$arcsin\left(\dfrac{\sqrt{3}}{3}\right)$}

    
    \item 如图,已知正三棱锥 ${ABCD},{BC} = {CD} = {BD} = \sqrt{3},{AB} = {AC} = {AD} = 2$ ,点 $P,Q$ 分别在棱 ${BC},{CD}$ 上(不包含端点)，则直线 ${AP},{BQ}$ 所成的角的取值范围是\blkda{$\left( {\dfrac{\pi }{3},\dfrac{\pi }{2}}\right\rbrack$}
    
    \begin{center}
    \includegraphics[width=0.3\textwidth]{bodyxb/src/bo_d4ojskjef24c73bc0dj0_6_241_910_382_394_0.jpg}
    \end{center}

    \begin{jieda}
        如下图所示:
        
        \begin{center}
        \includegraphics[width=0.3\textwidth]{bodyxb/src/bo_d4ojskjef24c73bc0dj0_6_242_1650_380_393_0.jpg}
        \end{center}

        设$AP$和$BQ$所成角为$\alpha$，当$BQ$和平面${PAO}$垂直时，$\alpha = \dfrac{\pi}{2}$，可取到.

        设$P$的线面角为$\angle{APO}$，当$P$点确定时，根据最小角定理可知$\alpha \geq \angle{APO}$（对于特定的$P$未必能取等，即$CD$上未必有$Q$使得$BQ // PO$）

        根据最大角定理可知$\angle{APO} \geq \angle{ABO} = \angle{ACO} = \dfrac{\pi}{3}$，当$P$取$B$，$Q$为$CD$中点时$\alpha = \dfrac{\pi}{3}$，故${AP},{BQ}$ 所成的角$\alpha$的下限(取不到)为$\dfrac{\pi}{3}$
        
        因为$P,Q$变化具有连续性，故${AP},{BQ}$ 所成的角的取值范围是$\left( {\dfrac{\pi }{3},\dfrac{\pi }{2}}\right\rbrack$.
    \end{jieda}
\end{enumerate}


\section*{12月18日}
\subsection*{动点与平行结合}
\begin{enumerate}

    \item 在棱长为 2 的正方体 ${ABCD} - {A}_{1}{B}_{1}{C}_{1}{D}_{1}$ 中，点 $E,F$ 分别是棱 ${BC},{C{C}_{1}}$ 的中点，动点 $P$ 是侧面 ${BC}{C}_{1}{B}_{1}$ 内一点，若 ${A}_{1}P//$ 平面 ${AEF}$ ，则点 $P$ 的轨迹长度为\blkda{$\sqrt{2}$}.
    \begin{center}
        \includegraphics[width=0.3\linewidth]{bodyxb//src/lifangti.png}
    \end{center}

    \begin{jieda}
        如下图所示,分别取棱 $B{B}_{1},{B}_{1}{C}_{1}$ 的中点 $M\text{ 、 }N$ ,连接 ${MN},B{C}_{1}$
        
        \begin{center}
        \includegraphics[width=0.3\textwidth]{bodyxb/src/bo_d50clvjef24c73bfi5dg_3_241_1779_393_349_0.jpg}
        \end{center}

        $MN // EF$，$A_1N // AE$，故平面$A_1MN // $平面$AEF$，故$P$的轨迹为$MN$，因此其长度为$\sqrt{2}$.
    \end{jieda}

    \item 已知棱长为 1 的正方体 ${ABCD} - {A}_{1}{B}_{1}{C}_{1}{D}_{1}$ ， $M$ 是 $B{B}_{1}$ 的中点，动点 $P$ 在正方体内部或表面上,且 ${MP}//$ 平面 ${AB}{D}_{1}$ ,则动点 $P$ 的轨迹所形成区域的面积是\blkda{$\dfrac{\sqrt{2}}{2}$}.
    \begin{center}
        \includegraphics[width=0.3\linewidth]{bodyxb//src/lifangti.png}
    \end{center}
    
    \begin{jieda}
        过点 $M$ 作平面 ${AB}{D}_{1}$ 的平行截面,再求解四边形的面积即可.

        \begin{center}
        \includegraphics[width=0.3\textwidth]{bodyxb/src/bo_d50clvjef24c73bfi5dg_5_243_199_492_470_0.jpg}
        \end{center}
    \end{jieda}
    
    \item 如图,在棱长为 2 的正方体 ${ABCD} - {A}_{1}{B}_{1}{C}_{1}{D}_{1}$ 中, $E,F,O$ 分别为正方形 ${AB}{B}_{1}{A}_{1}$，${A}_{1}{B}_{1}{C}_{1}{D}_{1}$ , ${ABCD}$ 的中心,点 $P$ 在正方形 ${ABCD}$ 内 (含边界) 运动,若直线 ${B}_{1}P$ 与平面 ${DEF}$ 无交点, 则点 $P$ 所形成的轨迹长度为\blkda{$\dfrac{2\sqrt{10}}{3}$}.
    
    \begin{center}
    \includegraphics[width=0.3\textwidth]{bodyxb/src/bo_d50clvjef24c73bfi5dg_0_244_716_446_394_0.jpg}
    \end{center}

    \begin{jieda}
        ${B}_{1}P$ 与平面 ${DEF}$ 无交点即${B}_{1}P$ 与平面 ${DEF}$ 平行，故动直线$B_1P$在一与平面 ${DEF}$ 平行的平面内

        $B_1O // DF$，取$CC_1D_1D$的中心$O'$，则$B_1O' // ED$，故平面$B_1OO' // $平面 ${DEF}$
        
        平面$B_1O'O$和平面$ABCD$的交线即为$P$的轨迹，易知$B_1O'$延长后与$AD$相交，设交点为$G$，设直线$OG$和$AB, CD$分别交于$M,N$，则$MN$即为$P$的轨迹，根据勾股定理可以求得其长度为$\dfrac{2\sqrt{10}}{3}$
    \end{jieda}

    \item 如图，在四棱锥 $P - {ABCD}$ 中， ${PA}\bot$ 平面 ${ABCD}$ ，底面 ${ABCD}$ 为正方形， ${PA} = {AB} = 4$ ， 点 $E,F$ 分别为 ${CD},{CP}$ 的中点,点 $T$ 为 $\bigtriangleup {PAB}$ 内的一个动点 (包括边界)， ${CT}//$ 平面 ${AEF}$ ，则点 $T$ 的轨迹的长度为\blkda{$\dfrac{2\sqrt{5}}{3}$}.
    
    \begin{center}
    \includegraphics[width=0.3\textwidth]{bodyxb/src/bo_d50clvjef24c73bfi5dg_2_241_611_346_340_0.jpg}
    \end{center}

    \begin{jieda}
        取$AB$中点$M$，则$CM // AE$，连接$BD$交$CM$于$N$，则$\dfrac{BN}{DN} = \dfrac{1}{2}$，取$BP$上靠近$B$的三等分点$Q$，则有$QN // PD // EF$，因此平面$CMQ //$平面$AEF$，故$T$的轨迹为$MQ$，根据余弦定理可求得其长度为$\dfrac{2\sqrt{5}}{3}$
    \end{jieda}
\end{enumerate}
\subsection*{动点与垂直结合}
\begin{enumerate}
    \item 如图，正四棱锥 $P - {ABCD}$ 的底面边长为 2,高为 2, $E$ 是边 ${BC}$ 的中点，动点 $Q$ 在表面上运动， 并且总保持 ${QE}\bot {AC}$ ,则动点 $P$ 的轨迹的周长为\blkda{$\sqrt{2} + \sqrt{6}$}.
    \begin{center}
        \includegraphics[width=0.3\linewidth]{bodyxb//src/dongdian1.png}
    \end{center}
    \begin{jieda}
        $Q$的轨迹为过$E$的$AC$的垂面在正四棱锥 $P - {ABCD}$上的截面

        取$PC$和$CD$的中点$F,G$，取$EG$中点$M$，则$FM // PO$，故$FM \bot $平面$ABCD$，故$FM  \bot AC$，且$EG \bot AC$，所以$AC \bot $平面$EFG$，故$P$的轨迹为$\triangle EFG$，其周长为$\sqrt{2}+\sqrt{6}$
    \end{jieda}
    
    \item 已知四棱柱 ${ABCD} - {A}_{1}{B}_{1}{C}_{1}{D}_{1}$ 的底面为菱形， $A{A}_{1} \bot$ 底面 ${ABCD},A{A}_{1} = 6,{AB} = 6$ ， $\angle {BCD} = {60}^{ \circ  }$ ,点 $M$ 是线段 ${BC}$ 上靠近 $C$ 的三等分点,点 $N$ 在四棱柱 ${ABCD} - {A}_{1}{B}_{1}{C}_{1}{D}_{1}$ 的表面上移动，且 ${MN}\bot {A}_{1}C$ ，则动点 $N$ 的轨迹长度为\blkda{$2 + 2\sqrt{13}$}.

    \begin{jieda}
        如图所示
            
        \begin{center}
        \includegraphics[width=0.3\textwidth]{bodyxb/src/bo_d50clvjef24c73bfi5dg_8_241_1196_358_343_0.jpg}
        \end{center}
    
        取$CD$上靠近$C$的三等分点，则$MF \bot AC$，取$CC_1$中点$G$，故平面$MFG \bot $平面$ACC_1A_1$，因为$AC \bot OG$，所以$AC \bot$平面$GMF$，故$N$的轨迹为$\triangle MFG$，则动点 $N$ 的轨迹长度为 ${MF} + {MG} + {FG} = 2 + 2\sqrt{13}$ .
    \end{jieda}
\end{enumerate}
\subsection*{动点与距离结合}
\begin{enumerate}
    \item 已知正三棱锥 $P - {ABC}$ 的六条棱长均为 $6,S$ 是 $\bigtriangleup {ABC}$ 及其内部的点构成的集合. 设集合 $T = \left\{  {Q \mid  {PQ} \leq  2\sqrt{7},Q \in  S}\right\}$ ,则 $T$ 表示的区域的面积为\blkda{${2\pi } + 3\sqrt{3}$}
    \begin{jieda}
        设顶点 $P$ 在底面上的投影为 $O$ ,连接 ${BO}$ ,则 $O$ 为三角形 ${ABC}$ 的中心,
        
        且 ${BO} = \dfrac{2}{3} \times  6 \times  \dfrac{\sqrt{3}}{2} = 2\sqrt{3}$ ,故 ${PO} = \sqrt{P{B}^{2} - B{O}^{2}} = \sqrt{{6}^{2} - {\left( 2\sqrt{3}\right) }^{2}} = 2\sqrt{6}$ .
        
        \begin{center}
        \includegraphics[width=0.2\textwidth]{bodyxb/src/bo_d50clvjef24c73bfi5dg_10_241_286_286_391_0.jpg}
        \end{center}
        
        
        因为 ${PQ} \leq  2\sqrt{7}$ ,故 ${OQ} = \sqrt{P{Q}^{2} - P{O}^{2}} \leq  \sqrt{{\left( 2\sqrt{7}\right) }^{2} - {\left( 2\sqrt{6}\right) }^{2}} = 2$ ,
        
        故 $Q$ 在以 $O$ 为圆心,2 为半径的圆及其内部,
        
        而三角形 ${ABC}$ 内切圆的圆心为 $O$ ,半径为 $\dfrac{1}{3} \times  6 \times  \dfrac{\sqrt{3}}{2} = \sqrt{3}$ ,而 $\sqrt{3} < 2 < 2\sqrt{3}$ ,
        
        故 $T$ 表示的区域为以 $O$ 为圆心,2 为半径的圆在 $\bigtriangleup  {ABC}$ 内(包含边界)的部分，
        
        设该圆与 ${AB}$ 交于 $D,E$ 两点,则 ${DE} = 2\sqrt{{2}^{2} - {\left( \sqrt{3}\right) }^{2}} = 2$ ,
        
        即 $\bigtriangleup  {ODE}$ 为正三角形,则弧 ${DE}$ 所在的弓形的面积为 $\dfrac{1}{2} \times  \dfrac{\pi }{3} \times  4 - \dfrac{\sqrt{3}}{4} \times  {2}^{2} = \dfrac{2\pi }{3} - \sqrt{3}$ ,
        
        故 $T$ 表示的区域的面积为 $\pi  \times  {2}^{2} - 3\left( {\dfrac{2\pi }{3} - \sqrt{3}}\right)  = {2\pi } + 3\sqrt{3}$
    \end{jieda}
\end{enumerate}
\subsection*{动点与角度结合}
\begin{enumerate}
    \item 已知平面 ${ABCD} \bot$ 平面 ${ADEF}$ ， ${AB} \bot  {AD},{CD} \bot  {AD}$ ，且 ${AB} = 3$ ， ${AD} = {CD} = 6$ ， ${ADEF}$ 是正方形，在正方形 ${ADEF}$ 内部(含边界)有一点 $M$ ，满足 ${MB}$ ， ${MC}$ 与平面 ${ADEF}$ 所成的角相等，则点 $M$ 的轨迹长度为\blkda{$\dfrac{4\pi }{3}$}.

    \begin{jieda}
        因为平面 ${ABCD} \bot$ 平面 ${ADEF}$ ,平面 ${ABCD} \cap$ 平面 ${ADEF} = {AD}$ ,
        
        ${AB} \bot  {AD},{CD} \bot  {AD}$ ,
        
        且 ${AB},{CD} \subset$ 平面 ${ABCD}$ ,所以 ${AB} \bot$ 平面 ${ADEF},{CD} \bot$ 平面 ${ADEF}$ ,
        
        可知 $\angle {BMA}$ 和 $\angle {CMD}$ 分别是直线 ${MB},{MC}$ 与平面 ${ADEF}$ 所成的角,
        
        则 $\angle {BMA} = \angle {CMD}$ ,即 $\tan \angle {BMA} = \tan \angle {CMD}$ ,
        
        可得 $\dfrac{AB}{AM} = \dfrac{CD}{MD}$ ,所以 ${MD} = {2AM}$ .

        故$M$的轨迹为阿氏圆（的一部分）

        以 $A$ 为原点建立平面直角坐标系如下图所示,
        \begin{center}
        \includegraphics[width=0.2\textwidth]{bodyxb/src/bo_d50clvjef24c73bfi5dg_9_569_488_343_285_0.jpg}
        \end{center}

        由 ${MD} = {2AM}$ 得 $\sqrt{{\left( x - 6\right) }^{2} + {y}^{2}} = 2\sqrt{{x}^{2} + {y}^{2}}$ ,化简得 ${\left( x + 2\right) }^{2} + {y}^{2} = {16}$
        
        由于点 $M$ 在第一象限内，所以 $M$ 点的轨迹是以 $G\left( {-2,0}\right)$ 为圆心，半径为 4 的圆在第一象限的部分,
        
        令 $x = 0$ 代入原方程,解得 $y =  \pm  2\sqrt{3}$ ,故 $H\left( {0,2\sqrt{3}}\right)$ ,
        
        由于 ${GA} = 2$ ，所以 $\angle {HGA} = \dfrac{\pi }{3}$ ，
        
        所以点 $M$ 的轨迹长度为 $\dfrac{\pi }{3} \times  4 = \dfrac{4\pi }{3}$ .
    \end{jieda}
\end{enumerate}
\section*{12月25日}
\subsection*{角的最值}
\subsubsection*{知识点}
\begin{enumerate}
    \item 最小角定理: 平面的斜线和它在平面内的射影所成的锐角 $\theta$ ,是这条斜线和平面内任一直线 ${AB}$ 所成角 $\angle {PAB}$ 中的最小值, 即线面角是最小的线线角. (如图, 由三余弦定理 $\cos \angle {PAB} = \cos \theta  \cdot  \cos \angle {OAB}$ 可得)
    
    \begin{center}
    \twotu{\includegraphics[width=0.45\textwidth]{bodyxb/src/bo_d1ig044601uc738hj6v0_68_1111_485_407_301_0.jpg}}{\includegraphics[width=0.45\textwidth]{bodyxb/src/bo_d1ig044601uc738hj6v0_68_1124_840_359_362_0.jpg}}
    
    \end{center}

    \item 最大角定理:对于一个锐二面角,在其中一个半平面内的任一条直线与另一个半平面所成的线面角 $\theta$ 的最大值等于二面角的平面角 $\alpha$ ,即二面角是最大的线面角. (如图，由三正弦定理 $\sin \theta  = \sin \alpha  \cdot  \sin \angle {PAB}$ 可得)
    
    \item 若二面角的一个半平面内有一条直线垂直于公共棱, 则此二面角的平面角等于该直线与另一个半平面所成的角, 即二面角与线面角可以相互转换.
\end{enumerate}
\subsubsection*{习题}
\begin{enumerate}
    \item 已知正三棱柱 ${ABC} - {A}^{\prime }{B}^{\prime }{C}^{\prime }$ 的所有棱长均为 2,且点 $P$ 在 ${B}^{\prime }{C}^{\prime }$ 上运动,则直线 ${AP}$ 与平面${ABC}$ 所成角的最大正弦值为\blkda{$\dfrac{2\sqrt{7}}{7}$}
    \begin{jieda}
        根据最大角定理可知且点 $P$ 在 ${B}^{\prime }{C}^{\prime }$ 上运动,则直线 ${AP}$ 与平面${ABC}$ 所成角的最大值为平面$B'C'A$和平面$BCA$所成锐二面角的的大小.
    \end{jieda}
    \item 如图,在四边形 ${ABCD}$ 中, ${AB} = {BC} = {CD} = 2,{AD} = \sqrt{5}$ ,对角线 ${BD} = 3,E$ 是直线 ${CD}$ 上任一点，将 $\bigtriangleup {ABD}$ 沿 ${BD}$ 翻折成 $\bigtriangleup {A}^{\prime }{BD}$ ,使二面角 ${A}^{\prime } - {BD} - C$ 的大小为 ${120}^{ \circ  }$ , 设异面直线 ${A}^{\prime }D$ 和 ${BE}$ 所成的角为 $\alpha$ ,则 $\sin \alpha$ 的最小值是\blkda{$\dfrac{\sqrt{3}}{3}$}.
    \begin{center}
            
        \twotu{\includegraphics[width=0.5\textwidth]{bodyxb/src/bo_d1ig044601uc738hj6v0_69_813_1076_298_264_0.jpg}}{\includegraphics[width=0.5\textwidth]{bodyxb/src/bo_d1ig044601uc738hj6v0_69_1138_1072_342_264_0.jpg}}
    \end{center}
    \begin{jieda}
        当$sin\alpha$最小的时候，$\alpha$最小，根据最小角定理可知此时${A}^{\prime }D$ 和 ${BE}$ 所成的角等于${A}^{\prime }D$ 和平面$BCD$所成角，设$A'$在面$BCD$内射影为$H$，故计算可得$sin\alpha = \dfrac{A'H}{A'D} = \dfrac{\dfrac{\sqrt{15}}{3}}{\sqrt{5}} = \dfrac{\sqrt{3}}{3}$
    \end{jieda}
    \clearpage
    \item 在四面体 ${ABCD}$ 中, $\bigtriangleup {BCD}$ 为等边三角形, $\angle {ADB} = \dfrac{\pi }{2}$ ,二面角 $B - {AD} - C$ 的大小为 $\alpha$ , 则 $\alpha$ 的取值范围是\blkda{$\left(0, \dfrac{\pi}{3}\right]$}
    \begin{center}
    \includegraphics[width=0.3\textwidth]{bodyxb/src/bo_d1ig044601uc738hj6v0_69_456_1053_312_290_0.jpg}
    \end{center}
    \begin{jieda}
        因为$DB \bot AD$，所以$DB$和平面$ACD$的夹角即为$B - {AD} - C$的平面角，根据最小角定理可知$\alpha \leq \angle{BDC} = \dfrac{\pi}{3}$
    \end{jieda}
    
    \item 如图,在三棱柱 ${ABC} - {A}_{1}{B}_{1}{C}_{1}$ 中, ${AB}\text{ 、 }{AC}\text{ 、 }A{A}_{1}$ 两两互相垂直, ${AB} = {AC} = A{A}_{1},M\text{ 、 }N$ 是线段 $B{B}_{1}\text{ 、 }C{C}_{1}$ 上的点,平面 ${AMN}$ 与平面 ${ABC}$ 所成(锐)二面角为 $\dfrac{\pi }{6}$ . 当 $\left| {{B}_{1}M}\right|$ 最小时, $\angle {AMB} =$ \blkda{$\dfrac{\pi }{3}$}
    
    \begin{center}
    \includegraphics[width=0.25\textwidth]{bodyxb/src/bo_d1ig044601uc738hj6v0_69_172_1006_243_337_0.jpg}
    \end{center}
    \begin{jieda}
        易知三棱柱 ${ABC} - {A}_{1}{B}_{1}{C}_{1}$为直三棱柱，$\left| {{B}_{1}M}\right|$ 最小即$\angle{AMB}$最小.
        
        根据最大角定理可知$AM$和平面$ABC$的线面角（$\angle{MAB}$）的最大值为平面 ${AMN}$ 与平面 ${ABC}$ 所成(锐)二面角，故$\angle{MAB} \leq \dfrac{\pi}{6}$，即$\angle{AMB} \geq \dfrac{\pi}{3}$
    \end{jieda}
    \item 设三棱锥 $V - {ABC}$ 的底面是正三角形,侧棱长均相等, $P$ 是棱 ${VA}$ 上的点 (不含端点). 记直线 ${PB}$ 与直线 ${AC}$ 所成的角为 $\alpha$ ,直线 ${PB}$ 与平面 ${ABC}$ 所成的角为 $\beta$ ,二面角 $P - {AC} - B$ 的平面角为 $\gamma$ ,则 \dotfill{(\kh{B})}
    \xx{$\beta  < \gamma ,\alpha  < \gamma$}{$\beta  < \alpha ,\beta  < \gamma$}{$\beta  < \alpha ,\gamma  < \alpha$}{$\alpha  < \beta ,\gamma  < \beta$}
    \begin{jieda}
        根据最小角定理可知$\alpha \geq \beta$，易知不满足取等条件，故$\beta < \alpha$

        \ps 因为$PB$并不在二面角$P - {AC} - B$的半平面内，因此不适用最大角定理

        过$P$做$PM \bot $面$ABC$，做$MN \bot AC$，则$\beta = \angle{PBM}$，$\gamma = \angle{PNM}$，故$tan\gamma = \dfrac{PM}{MN}$，$tan\beta = \dfrac{PM}{BM}$

        $V - {ABC}$为正四面体，根据正四面体的性质可知$M$在$\angle{BAC}$的中垂线上，有$MB = MC$，因为$MC > MN$，所以$tan\gamma > tan\beta$
    \end{jieda}
    
    \item 如图,在 $\bigtriangleup {ABC}$ 中, ${AB} = 1,{BC} = 2\sqrt{2},\angle B = \dfrac{\pi }{4}$ ,将 $\bigtriangleup {ABC}$ 绕边 ${AB}$ 翻转至 $\bigtriangleup {ABP}$ ,使平面 ${ABP} \bot$ 平面 ${ABC},D$ 是边 ${BC}$ 的中点,设 $Q$ 是线段 ${PA}$ 上的动点,当 ${PC}$ 与 ${DQ}$ 所成的角取得最小值时,线段 ${AQ}$ 的长度为\blkda{$\dfrac{2\sqrt{5}}{5}$}.
    
    \begin{center}
    \includegraphics[width=0.3\textwidth]{bodyxb/src/bo_d1ig044601uc738hj6v0_69_1169_124_331_299_0.jpg}
    \end{center}
    \begin{jieda}
        根据最小角定理可知，当 ${PC}$ 与 ${DQ}$ 所成的角取得最小值时，$DQ$与$PC$在面$APD$内射影平行

        延长$BA$至$O$使得$AO = 1$，则$\angle{BOC} = \dfrac{\pi}{2}$，故$CO \bot $平面$BOP$，且$PO \bot $平面$BOC$（$O$为墙角）

        将图形补成立方体，补全$PAD$对应截面，并做出$PC$在面内的投影，根据比例关系可以求得$AQ = \dfrac{2\sqrt{5}}{5}$
    \end{jieda}
    
\end{enumerate}